% 中文阅读版。正文（第 1--8 节）为中文译文；附录暂保留英文。与 ../p1.tex 同一版本（v0.9）。
\documentclass[11pt,UTF8]{ctexart}
\usepackage[margin=1in]{geometry}
\usepackage{amsmath,amssymb,amsthm,mathtools}
\usepackage{booktabs,array}
\usepackage{xcolor}
\usepackage[colorlinks=true,linkcolor=blue!50!black,citecolor=blue!50!black,urlcolor=blue!50!black]{hyperref}
\emergencystretch=3em

\newcommand{\unverified}[1]{\textcolor{red}{[#1 --- not yet checked against the original]}}
\newcommand{\draftnote}[1]{\textcolor{orange!70!black}{[Draft note: #1]}}

\newtheorem{theorem}{定理}[section]
\newtheorem{lemma}[theorem]{引理}
\newtheorem{proposition}[theorem]{命题}
\newtheorem{corollary}[theorem]{推论}
\theoremstyle{definition}
\newtheorem{definition}[theorem]{定义}
\newtheorem{example}[theorem]{例}
\theoremstyle{remark}
\newtheorem{remark}[theorem]{注}

\newcommand{\R}{\mathbb{R}}
\newcommand{\Z}{\mathcal{Z}}
\newcommand{\C}{\mathcal{C}}
\newcommand{\wfs}{\operatorname{wfs}}
\newcommand{\dist}{\operatorname{dist}}
\newcommand{\inter}{\operatorname{int}}
\newcommand{\cl}{\operatorname{cl}}
\newcommand{\conv}{\operatorname{conv}}

\title{容差 B-rep 何时定义一个实体？\\ 实现的存在性、不存在性与唯一性}
\author{张国庆（Eddie Zhang）\\ ZENFLUX, Miami Gardens, Florida, USA}
\date{中文阅读版，对应英文稿 v0.9，2026 年 11 月 13 日。仅供审阅，不外传。}

\begin{document}
\maketitle

\begin{abstract}
CAD 内核里的边界表示（B-rep）都带容差：边曲线、相邻面的边界曲线和顶点并不严丝合缝地对上，它们描述的实体只是``隐含''的。我们把这件事说清楚。容差 B-rep 的一个\emph{实现}，是它的胞腔复形的一个 PL 嵌入，并且每个胞腔都落在自己的容差区里。

\emph{唯一性。}设每个壳的亏格至少为 1，各胞腔的承载集可缩，每个壳有一个乘积领圈，领圈包含该壳的所有容差区，也包含一张参考曲面，且参考曲面到领圈的包含映射是同伦等价；各领圈两两不交。则所有实现的像彼此（拓扑意义下）环境同痕，它们围成的实体也环境同痕。领圈条件可由一个距离界推出：容差区到参考曲面的距离小于参考曲面的弱特征尺寸。对多面体参考曲面，我们用锥半径给出弱特征尺寸的一个显式下界，关键引理在 Lean 中机器验证。

\emph{存在性。}两种构造在给定假设下给出实现。第一种在边的法圆盘里用扇形、在顶点处用锥，是对正则复形的带条件定理，其假设中有两条是非局部的。第二种针对紧容差的尖锐边，把面所在曲面重新求交；它假设重新裁剪得到的候选曲面已经是嵌入的，并把它变成 PL 实现。

\emph{不存在性。}在明确的条件下，如果两张平面面片在所有边区、顶点区的``影子''之外横截相交，则不存在实现。

\emph{实验。}OpenCascade 的有效性检查判定不了上述任何一个问题。一次普通的缝合操作就能产生一个楔形体：它通过 OpenCascade 的两项检查，却没有实现。一根通过两项检查的厚壁管有两个实现，它们围成的实体甚至不同胚。
\end{abstract}

\section{引言}\label{sec:intro}

\subsection{问题}
CAD 内核把实体存成边界表示：一个胞腔复形（面、边、顶点），每个胞腔上挂着几何数据。
\begin{itemize}
\item 每个面带一张曲面上裁剪出来的面片。
\item 每条边带一条三维曲线，另外对每个相邻面各带一条参数域曲线（pcurve）。
\item 每个顶点带一个点。
\end{itemize}
在浮点运算下，经过布尔运算、倒圆角、偏置、缝合之类的操作后，这些数据互相对不上：pcurve 映到曲面上的像偏离三维边曲线，边曲线的端点偏离顶点。Jackson 的局部容差方法~\cite{Jackson95}（Parasolid 方法的公开描述）给每个顶点、边、面都配一个容差；只要偏差在容差之内，数据就算一致。例如 OpenCascade 要求 pcurve 的像落在三维曲线的边容差之内，边的端点落在顶点的顶点容差之内~\cite{OCCT}。

这样一个模型到底描述的是哪个实体？Qi 和 Shapiro~\cite{QiShapiro06b} 描述了通行的回答：如果在容差区内存在一个具有相同边界组合结构的精确实体，这个表示就是有效的。由此引出两个问题：
\begin{enumerate}
\item \emph{存在性。}是否存在一张具有给定胞腔结构、并遵守全部容差的嵌入曲面？
\item \emph{唯一性。}如果有好几张，它们围成的实体在环境同痕意义下是否都相同？
\end{enumerate}
下游算法（布尔、偏置、网格划分、分析）都默认这两个问题的答案是``是''。

\subsection{我们证明了什么}
我们把 (1) 中的嵌入称为\emph{实现}：胞腔复形 $K$ 的一个 PL 嵌入 $h$，对每个胞腔 $\sigma$ 满足 $h(\inter\sigma)\subseteq\Z(\sigma)$，其中 $\Z(\sigma)$ 是 $\sigma$ 的容差区（第~\ref{sec:setting}~节）。

\begin{itemize}
\item \textbf{唯一性}（定理~\ref{thm:uniqueness}）。
  \begin{itemize}
  \item \emph{假设：}各胞腔的承载集可缩；每个壳有一个乘积领圈，包含该壳的容差区和一张被承载的参考曲面，且参考曲面的包含映射是同伦等价；各领圈两两不交；每个壳的亏格至少为 1。
  \item \emph{结论：}任意两个实现的\emph{像}环境同痕，它们围成的实体也环境同痕。
  \item \emph{方法：}承载引理（一个逐胞腔的同伦论证），加上 Waldhausen 的定理：乘积 $F\times I$ 中的不可压缩曲面可以被推成水平的~\cite[Prop.~3.1]{Waldhausen68}。
  \end{itemize}
\item \textbf{领圈条件}（第~\ref{sec:collar}~节）。
  \begin{itemize}
  \item 领圈条件可以由参考曲面 $M$ 的弱特征尺寸 $\wfs(M)$ 以下的一个距离界推出（引理~\ref{lem:collar}）。其中一步用到 Moise 的定理：同胚可以用 PL 同胚逼近~\cite[Thm.~35.2]{Moise77}。
  \item 对多面体 $M$，我们证明 $\wfs(M)\ge\lambda(M)>0$，其中 $\lambda$ 由 $M$ 的锥半径定义（推论~\ref{cor:lambda}）。其中的锥引理在 Lean~4（Mathlib）中机器验证。
  \item 我们用八叉树搜索计算 $\wfs(M)$ 的下界，并说明了浮点意义上的限定。
  \end{itemize}
\item \textbf{存在性}（第~\ref{sec:existence}~节）。有两种构造：
  \begin{itemize}
  \item 在边的法圆盘里放扇形，在顶点处放锥。这是对正则复形的带条件定理（定理~\ref{thm:lemma0}），证明在附录~\ref{sec:proof51}。它的假设包括边和顶点上的局部条件、整体不相交条件（S1），以及两条非局部的拓扑条件（V1(b)、P0）。
  \item 对紧容差的尖锐边，把面所在曲面重新求交（命题~\ref{prop:B}、\ref{prop:Bprime}；曲面情形的证明在附录~\ref{sec:proofBprime}）。这两个命题假设重新裁剪得到的候选曲面已经嵌入并落在容差区内，然后把它变成一个 PL 实现。
  \end{itemize}
\item \textbf{不存在性}（定理~\ref{thm:obstruction}）。在明确的条件下，若两张平面面片在所有边区、顶点区的投影（``影子''）之外横截相交，则不存在实现。
\end{itemize}

\subsection{实验说明了什么}
在 OpenCascade~8.0.1 中（第~\ref{sec:experiments}~节）：
\begin{itemize}
\item \texttt{BRepCheck\_Analyzer} 不管容差多大都接受模型；源码里它的容差值检查只有声明、没有实现。
\item \texttt{BRepAlgoAPI\_Check} 会对可实现的模型误报，也会接受容差超过名义曲面弱特征尺寸的模型。其中有一根厚壁管，它有两个实现，围成的实体不同胚：其中一个在孔里打了一个三叶结。
\item 把一个面略有错位的楔形体缝合起来，得到的模型通过两项检查，但由定理~\ref{thm:obstruction}，它没有实现。定理的假设已用精确有理数运算验证。
\end{itemize}

\subsection{相关工作}

\subsubsection*{容差建模}
\begin{itemize}
\item Segal~\cite{Segal90} 给多面体模型的每个顶点、边、面都分配了容差。
\item Jackson~\cite{Jackson95} 给出了局部容差 B-rep 方法。
\item Milenkovic~\cite{Milenkovic88} 在平面上、对统一的 $\varepsilon$ 证明：规范化的数据在 $\varepsilon/2$ 以内有一个精确实现。
\item Sakkalis、Shen 和 Patrikalakis~\cite{SSP01} 研究区间实体：用盒子覆盖一个已知精确实体的边界，在同胚意义下讨论。Sakkalis、Peters 和 Bisceglio~\cite[Thm.~10]{SPB04} 证明：一张闭曲面若落在一张光滑闭曲面的法向管里，并且与每根法向纤维只交一次，则与之环境同痕。那里的管半径受曲率和一个双法向距离限制。
\item Qi 和 Shapiro~\cite{QiShapiro06} 发展了 $\varepsilon$-实体性，它不要求存在精确实体~\cite{QiShapiro06b}。
\item Shapiro~\cite{Shapiro08} 给出了容差区之并与胞腔复形同伦等价的条件。他在 \S5.2（第 178 页）指出，单有同伦条件是不够的，因为有效的边界表示还必须是一个可定向、能分隔空间的 2-闭链；把这些概念推广到容差区之并上``并不立即清楚''，不过同痕结果中用到的分隔概念可能有用。他还把``精确胞腔落在各自容差区内''看作强得多的要求。实现恰好满足的就是这个包含关系。
\end{itemize}

\subsubsection*{构造精确实体}
\begin{itemize}
\item Andersson、Stewart 和 Zidani~\cite{ASZ07} 以及 Stewart 和 Zidani~\cite{SZ08} 通过扰动（Whitney 延拓、超限插值）从不一致的裁剪面片数据构造实体。他们\emph{假设}扰动后的面片是嵌入的：
  \begin{itemize}
  \item ASZ 说相邻面片需要事后检查~\cite[pp.~819--820]{ASZ07}；
  \item SZ 作了类似假设~\cite[p.~188]{SZ08}，并推导了一个法向量界，用来排除多余的相交。
  \end{itemize}
\item 在更早的一篇论文~\cite{ASZ06} 中，同一批作者用度数论证，给出了单张定义在公共参数域上的复合参数曲面为单射的可检验条件。他们指出，分别裁剪的 B-rep 面还需要额外的内外信息~\cite[\S3.7]{ASZ06}。
\end{itemize}

\subsubsection*{同痕}
\begin{itemize}
\item Chazal 和 Lieutier~\cite{ChazalLieutier07}、Chazal、Cohen-Steiner 和 Lieutier~\cite{CCSL09dcg} 研究了弱特征尺寸，以及在它以下各偏置之间的同痕。后者还证明了一般紧集的偏置水平集在 Hausdorff 扰动下的同痕~\cite[Thms.~4.2,~4.3]{CCSL09cg}。
\item Chazal 和 Cohen-Steiner~\cite[Thm.~2.1]{CCS05} 证明：与 $S$ 同胚、落在 $S$ 的一个加厚层里、并把加厚层两侧隔开的曲面，与 $S$ 同痕。这与我们唯一性论证的拓扑核心很接近。
\item Andersson、Peters 和 Stewart~\cite[Thm.~1]{APS00} 证明：一族单参数的曲边三角面片，若每个时刻都拼接正确、每片都是单射且 Jacobi 矩阵非奇异、相邻面片不相切，则各时刻之间由环境同痕相连。他们的``不相切''与我们命题~\ref{prop:Bprime} 中的``不折叠''条款 T 很接近。两者设定不同：他们假设扰动后的复形在每个时刻都拼接正确；我们从面片彼此不相接的数据出发，去构造一个实现。
\end{itemize}

\subsection{哪些是新的，哪些不是}
对每个结果，我们说明：前人已知什么，我们新增什么，查到了什么程度。``我们没有找到''指的是本文所引文献及其背后的检索，不是穷尽的综述。

\subsubsection*{唯一性（定理~\ref{thm:uniqueness}）}
\begin{itemize}
\item \emph{已知。}拓扑核心是经典的。Waldhausen~\cite[Prop.~3.1]{Waldhausen68} 证明 $F\times I$ 中的不可压缩曲面可以推成水平的。Chazal 和 Cohen-Steiner~\cite[Thm.~2.1]{CCS05} 证明：与 $S$ 同胚、落在 $S$ 的加厚层里并隔开其两侧的曲面，与 $S$ 同痕。Sakkalis、Peters 和 Bisceglio~\cite[Thm.~10]{SPB04} 对落在光滑曲面法向管内、与每根纤维只交一次的曲面证明了同痕。
\item \emph{我们新增。}
  \begin{itemize}
  \item 针对容差 B-rep 的表述：每个胞腔各有容差区，实现要把每个胞腔留在自己的区里。
  \item 承载引理（引理~\ref{lem:carrier}）。在承载集可缩时，它使每个实现在领圈内与参考曲面同伦。再加上领圈假设``参考曲面的包含映射是同伦等价''，就得到 Waldhausen 定理所需的 $\pi_1$ 同构（引理~\ref{lem:separation}）。所以仍需要一个代替``分隔''的假设，即这个同伦等价条件，但只对唯一一张参考曲面要求，而不像~\cite{CCS05} 那样对每个实现都要求；也不需要像~\cite{SPB04} 那样要求实现与每根纤维只交一次。
  \item 由弱特征尺寸得到领圈（引理~\ref{lem:collar}），适用于 PL 参考曲面，允许尖锐边和三角化的小圆角。\cite{SPB04} 的结果是对光滑闭参考曲面陈述的，管半径受曲率和双法向距离限制。两者各有自己的适用范围：我们的结果需要 PL 参考曲面。
  \end{itemize}
\item \emph{局限。}
  \begin{itemize}
  \item 领圈条件是一个假设。只有当参考曲面是该壳的被承载 PL 嵌入、每个容差区都在 $N_d(M)$ 内且 $d<r<\wfs(M)$、不同壳的闭 $r$-邻域两两不交、亏格至少为 1 时，引理~\ref{lem:collar} 才能给出它。
  \item 不含亏格为 0 的壳。
  \item 同痕是拓扑意义下的，并且是像的同痕，不是嵌入映射本身的同痕。
  \end{itemize}
\end{itemize}

\subsubsection*{显式下界 $\wfs(M)\ge\lambda(M)$（推论~\ref{cor:lambda}）}
\begin{itemize}
\item \emph{已知。}对紧半代数集 $M$ 有 $\wfs(M)>0$。Fu~\cite[\S5.3]{Fu85} 证明紧半代数集的距离函数只有有限多个临界值。di Rocco、Eklund 和 Gäfvert~\cite[Thm.~2.5]{dREG22} 陈述了正性并归功于 Fu。Chazal 和 Lieutier~\cite[Prop.~3.6]{ChazalLieutier07} 对有界子解析开集陈述了临界值有限，说明半代数情形是 Fu 的结果，并指出推广的证明在他们更早的一篇论文里。Fu 用一组有限不等式来定义半代数集；他的论证（Tarski--Seidenberg、有限的光滑子流形分解、他的引理 4.3 和 Sard 定理）并不依赖这一限制。认为它覆盖多面体曲面，是我们自己的理解。
\item \emph{我们新增。}一个用锥半径表示的显式下界，其关键的锥引理在 Lean~4 中机器验证；$\lambda(M)>0$ 只有纸面证明，没有形式化。我们没有找到这样的下界，但这不说明它就是新的。
\end{itemize}

\subsubsection*{存在性（定理~\ref{thm:lemma0}，命题~\ref{prop:B}、\ref{prop:Bprime}）}
\begin{itemize}
\item \emph{已知。}
  \begin{itemize}
  \item Milenkovic~\cite{Milenkovic88}：平面上、统一容差下的存在性。
  \item Andersson、Stewart 和 Zidani~\cite{ASZ06}：公共参数域上单张复合面片为单射的可检验条件。
  \item \cite{ASZ07,SZ08}：从不一致数据构造实体，但假设扰动后的面片是嵌入的。
  \item Andersson、Peters 和 Stewart~\cite[Thm.~1]{APS00}：沿一族始终拼接正确的面片的环境同痕。他们的``不相切''与我们的不折叠条款 T 很接近。
  \end{itemize}
\item \emph{我们新增。}在每个胞腔各有容差区的三维 B-rep 上，给出整个胞腔复形存在实现的充分条件。它们分两类，性质不同。
  \begin{itemize}
  \item 定理~\ref{thm:lemma0} 从数据上的条件直接构造实现。它只对正则复形证明，并且是带条件的：其中两条假设是非局部的。
  \item 命题~\ref{prop:B}、\ref{prop:Bprime} 针对紧容差的尖锐边，假设一张重新裁剪得到的候选曲面已经嵌入并落在容差区内（在命题~\ref{prop:Bprime} 中即条款 C 和 B1--B3），然后证明它给出一个实现；对曲面情形这需要一个 PL 逼近。它们检验的是一个给定的候选，而不是推出嵌入性。命题~\ref{prop:B} 针对平面面和三度顶点；命题~\ref{prop:Bprime} 需要 $C^2$ 正则的面参数化，不能有极点或锥顶。在曲面模型上，这些条款只在采样点上用双精度检查，条款 T 的顶点部分没有检查。
  \end{itemize}
\item \emph{查到的程度。}对每个胞腔各有容差区的三维 B-rep，这两类存在性结果我们都没有找到。
\item \emph{未覆盖。}带相切边、并使用 OpenCascade 顶点容差的缝合模型。它们是否有实现，仍未解决（第~\ref{sec:open}~节）。
\end{itemize}

\subsubsection*{不存在性（定理~\ref{thm:obstruction}）}
\begin{itemize}
\item \emph{我们新增。}一个针对平面面片的可检验判据。我们把它用在通过 OpenCascade 两项检查的缝合模型上。
\item \emph{局限。}定理只涉及一对平面面，它们的闭胞腔是由单个线框围成的圆盘，且边区单连通（O3）。曲面版本仍未解决。
\item \emph{查到的程度。}在我们读过的文献里，没有找到容差 B-rep 的不存在性判据。
\end{itemize}

\subsubsection*{实验（第~\ref{sec:experiments}~节）}
\begin{itemize}
\item \emph{我们新增。}观察到 OpenCascade 的有效性检查既判定不了存在性，也判定不了唯一性。
\item \emph{范围。}证据是一个内核版本（8.0.1）里的几个零件。我们没有测试 ACIS 或 Parasolid。
\end{itemize}

\section{容差 B-rep 与实现}\label{sec:setting}

\subsection{复形}
$K$ 是 B-rep 意义下的有限胞腔复形。它的 0-胞腔是点，边是弧或闭曲线，面是紧曲面，可以带洞。开面是面去掉它所有的边（包括接缝）。底空间 $|K|$ 是一些连通闭曲面 $|K_1|,\dots,|K_n|$ 的不交并（称为\emph{壳}）。当 $\tau$ 是 $\sigma$ 的面（包括 $\tau=\sigma$）时记 $\tau\le\sigma$。

对唯一性定理，我们假设 $K$ 是\emph{正则}的：每个闭胞腔 $\bar\sigma$ 是一个球，其边界是子复形。带接缝边、闭边或带洞面的 B-rep 复形都不是正则的，需要先细分。新胞腔继承 $K$ 中包含它的那个开胞腔的容差区。这样容差区仍保持单调，而且原模型的实现也是细分后模型的实现，因为对新胞腔 $\sigma'\subseteq\sigma$ 有 $h(\inter\sigma')\subseteq h(\inter\sigma)\subseteq\Z(\sigma)$。此时第~\ref{sec:uniqueness}~节的承载条件要对细分后的复形成立。面 $f$ 内部的一条新边以 $\Z(f)$ 为承载集，所以细分面会增加假设 $\pi_1(\Z(f))=0$；可缩的容差区满足它。

$|K|$ 上的映射称为 \emph{PL} 的，如果它在 $|K|$ 的某个细分其胞腔的三角剖分的各单形上是线性的。在二维，任意两个这样的三角剖分都由一个保持胞腔的 PL 同胚相联系（标准事实；未单独对照原文），所以这个概念不依赖于剖分的选取。在边或面内部新产生的顶点继承该边或面的容差区，所以承载条件要求这些区道路连通。各存在性结果会分别说明它们在什么情况下直接适用于 B-rep 复形。

\subsection{几何数据与容差区}
\begin{itemize}
\item 每个顶点 $v$ 有一个点 $p_v$ 和容差 $\tau_v$。
\item 每条边 $e$ 有一条 $C^2$ 曲线 $\gamma_e\colon[0,1]\to\R^3$ 和容差 $\tau_e$。
\item 每个面 $f$ 有一张 $C^2$ 曲面 $S_f$、一块紧的裁剪面片 $P_f\subseteq S_f$ 和容差 $\tau_f$。$P_f$ 的边界由各边 $e<f$ 的 pcurve 的像 $\beta_{f,e}=S_f\circ(\text{pcurve})$ 组成。
\end{itemize}
边 $e$ 的\emph{偏差}是 $\mathrm{dev}_e=\max_{f>e}\max_t|\beta_{f,e}(t)-\gamma_e(t)|$，在相同参数处度量。OpenCascade 的 SameParameter 性质要求在这个意义下 $\mathrm{dev}_e\le\tau_e$；其格式规范把边容差定义为从 pcurve 的像到三维曲线的单侧距离~\cite{OCCT}。

\begin{definition}[容差区]\label{def:zones}
我们使用\emph{扩充的}开容差区：
\begin{itemize}
\item $\Z(v)=B(p_v,\tau_v)$；
\item $\Z(e)=N_{\tau_e}(\gamma_e)\cup\bigcup_{v<e}\Z(v)$；
\item $\Z(f)=N_{\tau_f}(P_f)\cup\bigcup_{e<f}\Z(e)$。
\end{itemize}
其中 $N_r$ 是开 $r$-邻域。扩充使 $\Z$ 单调：$\tau\le\sigma$ 蕴含 $\Z(\tau)\subseteq\Z(\sigma)$。
\end{definition}

\begin{definition}[实现]\label{def:realization}
容差 B-rep $T=(K,\text{数据},\tau)$ 的一个\emph{实现}是一个 PL 嵌入 $h\colon|K|\to\R^3$，对每个胞腔 $\sigma$ 满足 $h(\inter\sigma)\subseteq\Z(\sigma)$。
\end{definition}

\begin{remark}\label{rem:opencell}
\begin{itemize}
\item \emph{开胞腔条件。}我们只要求 $h(\inter\sigma)\subseteq\Z(\sigma)$。若用未扩充的邻域对闭胞腔提要求，即 $h(\bar\sigma)\subseteq N_{\tau_\sigma}(\cdot)$，通常对任何候选都不成立，主要因为面容差很小：$h(v)$ 将不得不落在每张相邻面片的 $\tau_f$ 之内。用扩充的容差区，开胞腔条件就蕴含 $h(\bar\sigma)\subseteq\Z(\sigma)$，因为 $\bar\sigma$ 是开胞腔 $\tau\le\sigma$ 的并。
\item \emph{驯顺性。}要求 $h$ 是 PL 的，排除了 Alexander 带角球面之类的野嵌入。
\end{itemize}
\end{remark}

\begin{definition}[承载集]
胞腔的\emph{承载集}是 $\C(\sigma)=\bigcup_{\tau\le\sigma}\Z(\tau)$。对扩充的容差区，$\C(\sigma)=\Z(\sigma)$。连续映射 $g\colon|K|\to\R^3$ 称为\emph{被承载的}，如果对每个 $\sigma$ 有 $g(\bar\sigma)\subseteq\C(\sigma)$。由注~\ref{rem:opencell}，每个实现都是被承载的。
\end{definition}

\begin{definition}[实体]\label{def:solid}
设 $S_1,\dots,S_n\subset\R^3$ 是两两不交的闭曲面。点 $x\notin\bigcup S_j$ 的\emph{深度}是使 $x$ 落在 $\R^3\setminus S_j$ 的有界分支内的 $j$ 的个数。由各 $S_j$ 围成的实体是深度为奇数的点集的闭包。
\end{definition}
当各壳定向一致时，这与 B-rep 中外壳与空腔的通常约定一致；定向一致性的检查另行处理。

\section{正亏格壳的唯一性}\label{sec:uniqueness}

\subsection{假设}
\begin{itemize}
\item \textbf{承载条件。}对每个顶点 $v$，$\C(v)$ 道路连通；对每个 $k\ge1$ 的 $k$-胞腔 $\sigma$，$\C(\sigma)$ 道路连通，且在每个基点处 $\pi_k(\C(\sigma))=0$。承载集可缩即可满足。
\item \textbf{领圈条件。}对每个壳 $K_i$，存在紧 PL 子流形 $U_i\subset\R^3$、闭曲面 $M_i$、PL 同胚 $\Psi_i\colon U_i\to M_i\times[-1,1]$，以及一个被承载的 PL 嵌入 $\varphi_i\colon|K_i|\to\R^3$（称为\emph{参考}），满足：
  \begin{itemize}
  \item $K_i$ 的所有胞腔的容差区都在 $\inter U_i$ 内；
  \item $\varphi_i(|K_i|)\subseteq\inter U_i$，且包含映射 $\varphi_i(|K_i|)\hookrightarrow U_i$ 是同伦等价；
  \item 各 $U_i$ 两两不交。
  \end{itemize}
  （参考曲面不必是某个水平层 $\Psi_i^{-1}(M_i\times\{t\})$。这个较弱的形式正是引理~\ref{lem:collar} 能给出的。）注意 $M_i$ 连通，因为 $M_i\simeq U_i\simeq|K_i|$；$M_i$ 可定向，因为 $M_i\times I\cong U_i\subset\R^3$；当 $|K_i|$ 不是球面时 $M_i$ 也不是球面，因为此时 $\pi_1(M_i)\cong\pi_1(|K_i|)\ne1$。
\item \textbf{亏格。}每个 $|K_i|$ 的亏格至少为 1。
\end{itemize}

\subsection{引理}

\begin{lemma}[承载引理]\label{lem:carrier}
在承载条件下，任意两个被承载的映射 $f,g\colon|K|\to\R^3$ 同伦，且同伦 $H$ 对每个 $\sigma$ 满足 $H(\bar\sigma\times I)\subseteq\C(\sigma)$。
\end{lemma}
\begin{proof}
对 $k$ 归纳，在 $|K|\times\{0,1\}\cup|K^{(k)}|\times I$ 上定义 $H$，保持 $H(\bar\sigma\times I)\subseteq\C(\sigma)$，且 $H(\cdot,0)=f$，$H(\cdot,1)=g$。
\begin{itemize}
\item $k=0$：对顶点 $v$，点 $f(v)$、$g(v)$ 都在道路连通集 $\C(v)$ 中；在 $\C(v)$ 中用一条道路连接它们。
\item $k\ge1$：设 $\sigma$ 是 $k$-胞腔。$\bar\sigma\times I$ 是一个 $(k+1)$-球，其边界 $\bar\sigma\times\{0,1\}\cup\partial\bar\sigma\times I$ 是一个 $k$-球面，$H$ 在它上面已有定义：
  \begin{itemize}
  \item 在 $\bar\sigma\times\{0,1\}$ 上是 $f$ 或 $g$，像在 $\C(\sigma)$ 中；
  \item 在 $\partial\bar\sigma\times I$ 上由归纳定义，像在 $\bigcup_{\tau<\sigma}\C(\tau)\subseteq\C(\sigma)$ 中。
  \end{itemize}
  由于 $\pi_k(\C(\sigma))=0$，$H$ 可以延拓到 $\bar\sigma\times I$ 上，像在 $\C(\sigma)$ 中。
\end{itemize}
各块在公共边界上一致，有限个闭集连续拼接得到连续的 $H$。
\end{proof}

\begin{lemma}[分隔]\label{lem:separation}
设承载条件和领圈条件成立，$h$ 是一个实现。则 $G=h(|K_i|)$ 是 $\inter U_i$ 中的闭曲面，且包含映射 $G\hookrightarrow U_i$ 是同伦等价。
\end{lemma}
\begin{proof}
\begin{itemize}
\item 对 $K_i$ 用引理~\ref{lem:carrier}，$h$ 与 $\varphi_i$ 通过像在 $\bigcup\C(\sigma)\subseteq\inter U_i$ 中的映射同伦，所以同伦发生在 $U_i$ 内。
\item $\varphi_i\colon|K_i|\to U_i$ 是到 $\varphi_i(|K_i|)$ 上的同胚再接包含映射；由领圈条件，后者是同伦等价。因此 $h\colon|K_i|\to U_i$ 也是同伦等价。
\item 由于 $h$ 是到 $G$ 上的同胚，包含映射 $G\hookrightarrow U_i$（等于 $h\circ h^{-1}$）是同伦等价。
\end{itemize}
\end{proof}

\begin{lemma}[水平化]\label{lem:horizontal}
设 $U\cong_{\Psi}M\times[-1,1]$ 如领圈条件中所述，$M$ 不是球面；$G\subset\inter U$ 是一张闭 PL 曲面，其到 $U$ 的包含映射是同伦等价。则存在 $\R^3$ 的环境同痕，支集在 $U$ 内、在 $\partial U$ 上不动，把 $G$ 变到 $\Psi^{-1}(M\times\{0\})$。
\end{lemma}
\begin{proof}
\begin{enumerate}
\item[(a)] \emph{$G$ 不可压缩。}
  \begin{itemize}
  \item 由假设，$\pi_1(G)\to\pi_1(U)$ 是同构。
  \item 压缩圆盘的边界会是一条在 $G$ 中本质、在 $U$ 中零伦的曲线，这被上一条排除。这就是 Waldhausen 意义下的不可压缩~\cite[p.~58]{Waldhausen68}。
  \item $G\simeq U\simeq M$，所以 $G$ 不是球面，该定义中的情形 (b) 不会出现。我们不用反方向（那需要环路定理）。
  \item $G$ 是可定向 3-流形中的可定向闭曲面，因而是双侧的，且与 $\partial U$ 不交；所以它满足 Waldhausen 的约定~\cite[\S0]{Waldhausen68}。
  \end{itemize}
\item[(b)] \emph{$G$ 可以推成水平的。}
  \begin{itemize}
  \item 对 $U\cong M\times I$（$M$ 可定向，不是球面）和不可压缩曲面 $G$ 应用 Waldhausen 的命题 3.1~\cite[pp.~65--68]{Waldhausen68}。由于 $\partial G=\emptyset$，条件``$\partial G\subseteq F\times1$''自动成立。
  \item 它给出一个在 $\partial U$ 上不动的同痕，把 $G$ 变成一张曲面 $G'$，其投影 $p\colon G'\to M$ 是到像上的同胚。
  \item $p(G')$ 在连通曲面 $M$ 中既开（区域不变性）又闭，所以 $p(G')=M$。因此 $G'$ 是某个函数 $s\colon M\to(-1,1)$ 的图像。
  \item Waldhausen 对命题 3.1 的证明用到 $F\times I$ 不可约：在基本情形（$F$ 为圆盘、环或裤子面，第 66 页）中明确使用，并通过第 67 页的归纳，用于切割得到的带边紧曲面 $F'$。下面的引理~\ref{lem:irreducible} 覆盖所有这些情形。
  \end{itemize}
\item[(c)] \emph{压平。}在每根纤维 $\{x\}\times[-1,1]$ 上取 PL 同胚，固定两端，把 $s(x)$ 送到 $(1-t)s(x)$，$t\in[0,1]$。它们给出 $U$ 的一个环境同痕，在 $\partial U$ 上不动，把 $G'$ 变到 $\Psi^{-1}(M\times\{0\})$。这个同痕是拓扑的。
\item[(d)] (b) 中的同痕本身就是环境同痕：在 Waldhausen 的论文里，$X$ 的子空间 $Z_1,Z_2$ 之间的同痕按定义是 $X$ 的恒等映射的一个同痕，把 $Z_1$ 变到 $Z_2$~\cite[p.~57]{Waldhausen68}；命题 3.1 给出的同痕在 $\partial U$ 上不动。两个同痕都在 $\partial U$ 上不动，可以用恒等映射延拓到 $\R^3$。
\end{enumerate}
\end{proof}

\begin{lemma}\label{lem:irreducible}
设 $F$ 是除 2-球面以外的紧可定向 PL 曲面。则 $F\times I$ 不可约：每个 PL 2-球面 $S\subset\inter(F\times I)$ 在 $F\times I$ 中围出一个 PL 3-胞腔。（按 Waldhausen 的约定，曲面是真嵌入的，所以 $F\times I$ 中的 2-球面与边界不交。）
\end{lemma}
\begin{proof}
\emph{情形 $\partial F\ne\emptyset$。}
\begin{itemize}
\item 此时 $F\times I$ 是一个柄体：$\R^3$ 中一个有限图的正则邻域 $H$，且 $\partial H$ 连通（标准事实；未单独对照原文）。我们在 $H\subset\R^3$ 中讨论。
\item 由 PL Schoenflies 定理~\cite[Thm.~17.12]{Moise77}，$S$ 在 $\R^3$ 中围出一个 PL 3-胞腔 $B$。集合 $\R^3\setminus H$ 连通（Alexander 对偶，因为 $\partial H$ 连通）、无界、与 $S$ 不交，所以它落在 $\R^3\setminus S$ 的无界分支里。因此 $B\subseteq H$。
\end{itemize}
\emph{情形 $F$ 闭、亏格至少为 1。}
\begin{itemize}
\item 设 $q\colon\tilde F\times I\to F\times I$ 是万有覆叠，甲板变换群 $\Gamma$ 以 $\gamma\times\mathrm{id}$ 作用。$\tilde F$ 是单连通、非紧、无边的 PL 曲面，所以 PL 同胚于 $\R^2$（标准事实；未单独对照原文）。因此 $\tilde F\times I\cong_{PL}\R^2\times I\subset\R^3$。
\item $S$ 单连通，所以它提升为一个 PL 2-球面 $\tilde S\subset\R^2\times(0,1)$，各提升 $\gamma\tilde S$（$\gamma\in\Gamma$）两两不交。由 PL Schoenflies 定理~\cite[Thm.~17.12]{Moise77}，$\tilde S$ 围出一个 PL 3-胞腔 $\tilde B\subset\R^3$。$\R^2\times[0,1]$ 之外的点在 $\tilde S$ 的补集中可以连到无穷远，所以 $\tilde B\subset\R^2\times I$。
\item 设 $\gamma\ne1$。球面 $\tilde S$ 与 $\gamma\tilde S$ 不交，所以胞腔 $\tilde B$ 与 $\gamma\tilde B$ 要么不交，要么一个套在另一个里。若 $\gamma\tilde B\subseteq\tilde B$，则对所有 $n\ge1$ 有 $\gamma^n\tilde B\subseteq\tilde B$。$\Gamma$ 无挠，所以各 $\gamma^n$ 互不相同，于是有无穷多个甲板变换使紧集 $\tilde B$ 与自身相交，与真不连续性矛盾。$\tilde B\subseteq\gamma\tilde B$ 的情形对 $\gamma^{-1}$ 同理。
\item 所以 $\tilde B$ 的各平移两两不交，$q$ 在 $\tilde B$ 上是单射，$q(\tilde B)$ 是 $F\times I$ 中以 $S$ 为边界的 PL 3-胞腔。
\end{itemize}
另外用到、未单独对照原文的标准事实：闭曲面的 $\pi_1$ 无挠；甲板变换群的作用是真不连续的。
\end{proof}

\subsection{定理}

\begin{theorem}[唯一性]\label{thm:uniqueness}
设承载条件、领圈条件成立，且每个壳的亏格至少为 1。则：
\begin{enumerate}
\item 任一实现的像都与 $\bigcup_i\Psi_i^{-1}(M_i\times\{0\})$ 环境同痕（通过一个拓扑同痕）；
\item 任意两个实现的像环境同痕；
\item 它们围成的实体（定义~\ref{def:solid}）环境同痕。
\end{enumerate}
\end{theorem}
\begin{proof}
\begin{itemize}
\item 由引理~\ref{lem:separation}，每个 $G_i=h(|K_i|)$ 在 $U_i$ 中满足引理~\ref{lem:horizontal} 的假设。逐个壳应用该引理。各支集 $U_i$ 两两不交，所以各同痕 $\Phi^i$ 可以同时进行，其终端映射是 $\R^3$ 的一个同胚 $\Phi$。
\item 对每个 $j$ 有 $\Phi(G_j)=\Psi_j^{-1}(M_j\times\{0\})$：在 $U_j$ 上，每个 $i\ne j$ 的 $\Phi^i$ 都是恒等映射，因为其支集 $U_i$ 与 $U_j$ 不交。
\item 由 Jordan--Brouwer 分隔定理，每张闭曲面把 $\R^3$ 分成一个有界分支和一个无界分支。$\R^3$ 的同胚把闭包紧的集合映成闭包紧的集合，所以 $\Phi$ 把 $\R^3\setminus G_j$ 的有界分支映成 $\R^3\setminus\Phi(G_j)$ 的有界分支。于是深度保持不变，$\Phi$ 把各 $G_j$ 围成的实体映成各水平层围成的实体。
\item 对两个实现及其映射 $\Phi$、$\Phi'$，取 $\Phi'^{-1}\circ\Phi$。
\end{itemize}
\end{proof}

\begin{remark}
\begin{itemize}
\item \emph{未覆盖的情形。}不含球面（亏格 0）；同痕是像的同痕，不是嵌入映射的同痕。
\item \emph{与 Chazal--Cohen-Steiner 的关系。}Chazal 和 Cohen-Steiner~\cite[Thm.~2.1]{CCS05} 证明：落在 $S$ 的加厚层里、与 $S$ 同胚并隔开加厚层两侧的紧可定向曲面 $S'$，在加厚层中与 $S$ 同痕。引理~\ref{lem:separation} 和~\ref{lem:horizontal} 在这里起同样的作用。容差 B-rep 特有的是整个设定：每个胞腔各有容差区、保持胞腔的实现、承载引理，以及由承载引理提供、用来代替``分隔''的 $\pi_1$ 条件。
\item \emph{领圈条件不能去掉。}第~\ref{sec:experiments}~节的厚壁管通过 OpenCascade 的两项检查，却有两个实现，围成的实体甚至不同胚。
\item \emph{用到而未单独对照原文的标准事实：}逐胞腔延拓的连续性、$\R^3$ 中闭曲面的可定向性、区域不变性、Jordan--Brouwer 分隔定理，以及按纤维压平是一个环境同痕。
\end{itemize}
\end{remark}

\section{从弱特征尺寸到领圈条件}\label{sec:collar}

\subsection{弱特征尺寸}
设 $M\subset\R^3$ 紧，$d(p)=\dist(p,M)$，$\Gamma(p)$ 是 $p$ 在 $M$ 中最近点的集合。点 $p\notin M$ 称为\emph{临界点}，如果 $0\in\conv\{p-x:x\in\Gamma(p)\}$。\emph{弱特征尺寸} $\wfs(M)$ 是 $d$ 的正临界值的下确界~\cite[\S2.2]{CCSL09dcg}；另见~\cite{ChazalLieutier07}。与 reach（可达距离）不同，它在尖锐边和小圆角处仍然是正的：带凹边的曲面 reach 为 $0$，半径 $r$ 的圆角把 reach 限制在 $r$ 以下。

\begin{lemma}[领圈]\label{lem:collar}
设 $M$ 是 $\R^3$ 中的连通闭 PL 曲面，$0<d<r<\wfs(M)$。
\begin{enumerate}
\item[(a)] 闭邻域 $\bar N_r(M)$ 同胚于 $M\times[-1,1]$，且 $M$ 对应 $M\times\{0\}$。
\item[(a$'$)] 存在紧 PL 子流形 $U$，满足 $\bar N_d(M)\subseteq\inter U$、$U\subseteq\inter\bar N_r(M)$ 且 $U\cong_{PL}M\times[-1,1]$。
\item[(b)] 设 $M$ 的亏格至少为 1，$U$ 如 (a$'$)。再设 $M=\varphi(|K_i|)$，其中 $\varphi$ 是某个壳的被承载 PL 嵌入；该壳每个胞腔的容差区都在 $N_d(M)$ 内；不同壳的 $\bar N_r(M_i)$ 两两不交。则第~\ref{sec:uniqueness}~节的领圈条件成立，取 $U_i=U$。
\end{enumerate}
\end{lemma}
\begin{proof}
记 $d=d_M$。由 Jordan--Brouwer 分隔定理，$\R^3\setminus M$ 有两个分支：有界的 $\Omega_-$ 和无界的 $\Omega_+$。

(a) \emph{第 1 步：小半径。}
\begin{itemize}
\item 在一个包含 $\bar N_r(M)$ 的大单纯立方体 $Q$ 中讨论，三角剖分使 $M$ 是子复形。$d$ 限制在 $Q$ 上是半代数的（Tarski--Seidenberg）、逆紧的、非负的，零点集是紧集 $M$。Durfee 对（半代数）\emph{rug 函数}的定义~\cite[Def.~3.1]{Durfee83} 要求的正是这些，不要求光滑性，所以 $d|_Q$ 是 $M$ 在 $Q$ 中的 rug 函数。于是对充分小的 $\delta>0$，集合 $\bar N_\delta(M)=d^{-1}[0,\delta]$ 是 $M$ 在 $Q$ 中的半代数邻域~\cite[Def.~3.4]{Durfee83}。
\item $M$ 在 $Q$ 中的 PL 正则邻域 $R$ 也是半代数邻域~\cite[Prop.~3.7]{Durfee83}。由半代数邻域的唯一性~\cite[Prop.~3.5]{Durfee83}，存在同胚 $h_t\colon Q\to Q$，对所有 $t$ 有 $h_t|_M=\mathrm{id}$，且 $h_1(\bar N_\delta(M))=R$。所以 $h_1$ 是一个逐点固定 $M$ 的偶同胚 $(\bar N_\delta(M),M)\to(R,M)$。取 $\delta<r$。
\item $M$ 是 $\R^3$ 中的可定向闭曲面，因而双侧，所以 $R\cong M\times[-1,1]$，$M$ 对应 $M\times\{0\}$。因此 $(\bar N_\delta(M),M)\cong(M\times[-1,1],M\times\{0\})$。它的边界是 $S_\delta=\{d=\delta\}$，因为 $\delta$ 不是临界值，每个 $d=\delta$ 的点都是 $d>\delta$ 的点的极限（第 3 步中的方向导数论证）。$S_\delta$ 有两个分支，即 $M\times\{\pm1\}$ 的像，它们位于 $\R^3\setminus M$ 的不同分支中：两半 $M\times(0,1]$ 和 $M\times[-1,0)$ 各自连通，所以各在某个 $\Omega_\pm$ 中；$M=\partial\Omega_+=\partial\Omega_-$ 上任一点的每个邻域都与两者相交，所以两半落在不同分支里。记 $S_\delta^\pm=S_\delta\cap\Omega_\pm$。
\end{itemize}

\emph{第 2 步：两侧的壳层。}令 $A_\pm=\{x\in\Omega_\pm:\delta\le d(x)\le r\}$。
\begin{itemize}
\item \emph{内侧。}取 $O=\Omega_-$。它有界且开；对 $x\in O$，到 $O^c$ 的距离等于 $d(x)$，最近点也相同（都在 $\partial O=M$ 上）。所以 $O$ 中到 $O^c$ 的距离的临界点，正是 $\Omega_-$ 中 $d$ 的临界点。两种临界点概念一致：\cite{ChazalLieutier07} 用的是其梯度场为零，而这当且仅当 $x$ 落在其最近边界点的凸包中~\cite[Lemma~2.2]{ChazalLieutier07}（另见~\cite[\S2.2]{CCSL09dcg}）。
\item 由于 $r<\wfs(M)$，壳层 $A_-$ 中没有临界点。同痕引理~\cite[Prop.~3.4]{ChazalLieutier07}（他们转述自 Grove 的距离函数临界点理论）给出同胚 $A_-\cong S^-_\delta\times[\delta,r]$。
\item \emph{外侧。}取 $O=\Omega_+\cap B(0,\mathcal R)$，$\mathcal R$ 很大。对 $x\in O$，到 $O^c$ 的距离是 $\min(d(x),\mathcal R-|x|)$。取 $\mathcal R$ 大于 $2r$ 加上 $|y|$（$y\in M$）的最大值。于是 $O$ 中到 $O^c$ 距离在 $[\delta,r]$ 内的点分成不交的两块：$A_+$，其中严格有 $d<\mathcal R-|x|$，所以距离等于 $d$ 且最近点相同；以及球面 $|x|=\mathcal R$ 附近的一个壳层，那里距离等于 $\mathcal R-|x|$。
\item 两块都不含临界点。在 $A_+$ 中是因为 $r<\wfs$；在球面附近，唯一的最近点是径向投影。
\item 于是命题 3.4 给出这两块之并与 $R^{-1}(\delta)\times[\delta,r]$ 的同胚，其中 $R^{-1}(\delta)=S^+_\delta\sqcup\{|x|=\mathcal R-\delta\}$ 恰有两个分支（$S^+_\delta\cong M$ 连通）。乘积的分支是 $C\times[\delta,r]$，$C$ 取这两个分支。球面附近的壳层是并集的一个分支，所以 $A_+$ 是另一个：$A_+\cong C\times[\delta,r]$，$C$ 连通。第 3 步说明 $C\cong S^+_\delta$。
\item 我们使用命题 3.4 原文所述的闭壳层形式：对有界开集 $O$ 和 $R=d(\cdot,O^c)$，若壳层 $\{x\in O:\rho_1\le R(x)\le\rho_2\}$ 不含 $R$ 的临界点，则它同胚于 $R^{-1}(\rho_1)\times[\rho_1,\rho_2]$~\cite[Prop.~3.4]{ChazalLieutier07}。
\end{itemize}

\emph{第 3 步：粘合。}
\begin{itemize}
\item \emph{$A_\pm$ 的边界。}$A_\pm$ 是 $\R^3$ 中的紧 3-流形，由区域不变性，它的流形边界就是它在 $\R^3$ 中的边界。
  \begin{itemize}
  \item $\delta<d<r$ 的点是内点。
  \item $d=\delta$ 的点是 $d<\delta$ 的点的极限：沿到最近点的线段移动即可。
  \item $d=r$ 的点 $x$ 是 $d>r$ 的点的极限。由于 $x$ 不是临界点，存在单位向量 $u$，对所有 $y\in\Gamma(x)$ 有 $\langle u,x-y\rangle>0$。$d$ 沿 $u$ 的单侧方向导数是 $\min_{y\in\Gamma(x)}\langle u,(x-y)/|x-y|\rangle>0$（距离函数方向导数的标准公式；未单独对照原文）。
  \end{itemize}
  所以 $\partial A_\pm=S^\pm_\delta\sqcup S^\pm_r$，连通闭集 $S^\pm_\delta$ 是一个边界分支。$C\times[\delta,r]$ 的每个边界分支都是 $C$ 的一个拷贝，所以 $C\cong S^\pm_\delta$。
\item 带边流形之间的同胚把边界分支映成边界分支。所以第 2 步的同胚把 $A_\pm$ 的边界分支 $S^\pm_\delta$ 映到乘积的某一端。必要时复合一个 $[\delta,r]$ 的翻转，再复合它在 $S^\pm_\delta$ 上的限制之逆，就得到 $A_\pm\cong S^\pm_\delta\times[0,1]$，$S^\pm_\delta$ 对应 $S^\pm_\delta\times\{0\}$。
\item 于是 $\bar N_r(M)=\bar N_\delta(M)\cup A_+\cup A_-$ 是 $M\times[-1,1]$ 沿每个边界分支各粘上一个乘积领圈。它同胚于 $M\times[-2,2]\cong M\times[-1,1]$，$M$ 仍对应 $M\times\{0\}$。
\end{itemize}

(a$'$) $\bar N_r(M)$ 的内部是 $N_r(M)$：由第 3 步的计算，$d=r$ 的点是 $d>r$ 的点的极限。
\begin{itemize}
\item 由 (a)，存在同胚 $g\colon M\times(-1,1)\to N_r(M)$。
\item $g^{-1}(\bar N_d(M))$ 紧，所以落在某个 $M\times[-a,a]$（$a<1$）中。取 $a<b<1$，令 $K=M\times[-b,b]$，它是 PL 3-流形 $M\times(-1,1)$ 中的带边多面体 3-流形。
\item 紧集 $g(K)$ 与闭集 $\R^3\setminus N_r(M)$ 不交，紧集 $g(\partial K)$ 与 $\bar N_d(M)$ 也不交。取 $\varepsilon>0$ 小于这两个距离。由 Moise 的逼近定理~\cite[Thm.~35.2]{Moise77}，存在 $K$ 到 $\R^3$ 中的 PL 同胚 $f$，在 $K$ 上 $|f-g|<\varepsilon$。令 $U=f(K)$。
\item 于是 $U$ 是紧 PL 子流形，$U\cong_{PL}M\times[-1,1]$，且 $U\subseteq N_r(M)=\inter\bar N_r(M)$。
\item $\bar N_d(M)\subseteq\inter U$。设 $x\in\bar N_d(M)$。从 $g|_{\partial K}$ 到 $f|_{\partial K}$ 的直线同伦避开 $x$，所以对两张边界曲面中的每一张，$x$ 是否落在它的有界补分支里，对 $f$ 和对 $g$ 结果相同（从 $x$ 出发的径向投影的模 2 度是同伦不变量）。对 $\R^3$ 中有两张边界曲面的紧连通 3-流形，其内部就是相对这两张曲面深度为 1 的点集（定义~\ref{def:solid}；标准事实，由 Alexander 对偶）。由于 $x\in g(\inter K)$，它对 $g$ 的深度为 1，因而对 $f$ 也为 1，所以 $x\in\inter U$。
\item $f$ 不一定把 $M$ 送到某个水平层，这就是第~\ref{sec:uniqueness}~节的领圈条件不作此要求的原因。
\end{itemize}

(b) 容差区在 $N_d(M)\subseteq\inter U$ 内，且 $M\subseteq N_d(M)\subseteq\inter U$。各 $U_i$ 两两不交，因为它们分别落在不交的 $\bar N_r(M_i)$ 中。剩下要证包含映射 $M\hookrightarrow U$ 是同伦等价。
\begin{itemize}
\item 由 (a)，复合 $M\hookrightarrow U\hookrightarrow\bar N_r(M)$ 是同伦等价，所以 $\pi_1(U)\to\pi_1(\bar N_r(M))$ 是满射。
\item 两个群都同构于 $\pi_1(M)$，即亏格至少为 1 的曲面群。曲面群是 Hopf 群，所以从它到一个同构拷贝的满射是同构。因此 $\pi_1(M)\to\pi_1(U)$ 是同构。
\item $M$ 和 $U\simeq M$ 都是非球面的（aspherical），由 Whitehead 定理，包含映射是同伦等价。
\end{itemize}
这里既用到 (a) 也用到 (a$'$)，并且用到亏格假设：对球面，论证不成立。用到而未单独对照原文的标准事实：Jordan--Brouwer 分隔定理、Tarski--Seidenberg、正则邻域理论（双侧闭曲面的正则邻域是乘积）、拓扑流形的边界不变性、沿边界分支粘乘积领圈、曲面群是 Hopf 群（它们剩余有限）、正亏格闭曲面的非球面性、Whitehead 定理、模 2 度的同伦不变性及其用于判断点在嵌入闭曲面哪一侧，以及两张不交闭曲面补集分支的 Alexander 对偶。
\end{proof}

\subsection{多面体曲面的显式下界}
对 $a\in M$，\emph{锥半径} $\rho(a)$ 是满足下述条件的 $r>0$ 的上确界：存在以 $a$ 为顶点的闭锥 $C$，使 $M\cap B(a,r)=C\cap B(a,r)$。

\begin{lemma}[锥引理]\label{lem:cone}
设 $a\in M$，$r<\rho(a)$，$p\notin M$ 且 $|p-a|<r/2$。则每个最近点 $x\in\Gamma(p)$ 满足 $\langle p-x,\,p-a\rangle=|p-x|^2$。特别地，$p$ 不是临界点。
\end{lemma}
\begin{proof}
\begin{itemize}
\item \emph{最近点落在球内。}$d(p)\le|p-a|<r/2$，所以每个 $x\in\Gamma(p)$ 满足 $|x-a|<r$。
\item \emph{过 $x$ 的射线留在 $M$ 中。}点 $y_s=a+s(x-a)$（$s$ 接近 $1$）落在锥内且在 $B(a,r)$ 内，因而在 $M$ 中。所以 $s\mapsto|p-y_s|^2$ 在 $s=1$ 处取局部极小，得 $\langle p-x,x-a\rangle=0$。
\item \emph{结论。}对所有 $x\in\Gamma(p)$，$\langle p-x,p-a\rangle=|p-x|^2+\langle p-x,x-a\rangle=d(p)^2>0$。所以 $0$ 不在向量 $p-x$ 的凸包中。
\end{itemize}
\end{proof}

\begin{corollary}\label{cor:lambda}
令 $\lambda(M)=\inf_{x\in M}\sup_{a\in M}\bigl(\rho(a)/2-|x-a|\bigr)$。则 $d$ 没有满足 $0<d(p)<\lambda(M)$ 的临界点 $p$，所以 $\wfs(M)\ge\lambda(M)$。若 $M$ 是有限个三角形之并，则 $\lambda(M)>0$。
\end{corollary}
\begin{proof}
\begin{itemize}
\item \emph{下界。}给定 $d(p)<\lambda$ 的 $p$ 和 $x\in\Gamma(p)$，取 $a$ 使 $\rho(a)/2-|x-a|>d(p)$。则 $|p-a|<\rho(a)/2$，引理~\ref{lem:cone} 适用。
\item \emph{正性。}对多面体，由多面体的局部锥结构，每个 $a$ 都有 $\rho(a)>0$（标准事实；未单独对照原文）。函数 $g(x)=\sup_a(\rho(a)/2-|x-a|)$ 是 1-Lipschitz 的且有限（因为 $\rho\le\operatorname{diam}M$），且 $g(x)\ge\rho(x)/2>0$。$M$ 紧，所以 $\min g>0$。
\end{itemize}
\end{proof}

引理~\ref{lem:cone} 及推论~\ref{cor:lambda} 的逐点形式，在 Lean~4（Mathlib）中对任意实内积空间的子集做了机器验证（附录~\ref{sec:lean}）。

\subsection{$\wfs$ 的认证下界}
对三角化的参考曲面 $M_{PL}$，我们用最佳优先八叉树计算一个 $L$，使得``$(c\cdot\mathrm{tol},L)$ 内没有临界值''。我们说 $L$ 是\emph{认证的}，仅限于以下限定意义：计算用双精度；判断点是否属于某个特征时用容差 $\mathrm{tol}=10^{-10}$ 乘以模型尺寸；搜索只在截断距离 $R$ 以下进行。
\begin{itemize}
\item \emph{靠近 $M$ 的盒子}用引理~\ref{lem:cone} 认证，用到候选点 $a$ 处 $\rho(a)$ 的一个下界。这个下界取 $a$ 到不包含 $a$ 的最近特征（平面区域、边界线段、顶点）的距离。对其有向边界线段覆盖了边界的平面区域，在对每个合并区域做过全局共面复查之后，这个下界是可靠的。
\item \emph{离得较远的盒子}用方向检验认证：一个小线性规划寻找单位向量 $u$，使盒内所有点 $p$ 的所有候选最近点 $x$ 都满足 $\langle u,p-x\rangle>0$。候选点是盒子中心到相关特征的投影。
\item \emph{停止规则。}搜索在遇到第一个太小而不能再分、又无法认证的盒子时停止。它到 $M$ 的距离就是认证下界 $L$。
\end{itemize}
结果见表~\ref{tab:wfs}。引理~\ref{lem:collar} 需要 $(0,r]$ 内没有临界值。计算得到的下界覆盖 $(c\cdot\mathrm{tol},L)$；当 $\lambda(M)>c\cdot\mathrm{tol}$ 时，区间 $(0,c\cdot\mathrm{tol}]$ 由推论~\ref{cor:lambda} 覆盖，这一点我们没有单独检查。

\begin{table}[h]
\centering
\begin{tabular}{lcc}
\toprule
零件 & 精确 $\wfs$（光滑零件） & $M_{PL}$ 的认证下界 \\
\midrule
立方体，边长 10 & 5 & 4.530 \\
板 $4\times4\times1$ & 0.5 & 0.436 \\
板 $4\times4\times0.3$ & 0.15 & 0.137 \\
L 形（壁厚 2） & 1 & 0.894 \\
L 形，内圆角 0.2 & 1 & 0.894 \\
带 0.2 宽槽的块 & 0.1 & 0.086 \\
圆柱 $R=3$，$H=10$ & 3 & 2.512 \\
\bottomrule
\end{tabular}
\caption{三角化零件的 $\wfs$ 下界，在上述限定意义下认证。``精确''一列是光滑零件的 $\wfs$；对三角化的圆柱，$\wfs(M_{PL})$ 略小于 $3$。内圆角把 reach 限制在 $0.2$，但下界不变。}\label{tab:wfs}
\end{table}

\section{存在性}\label{sec:existence}

我们给出两种构造，把数据上的假设变成一个实现。两者都不假设扰动后的曲面整体嵌入。定理~\ref{thm:lemma0} 确实假设（P0）每张原始面片是嵌入的，且其假设中有两条是非局部的。命题~\ref{prop:B}、\ref{prop:Bprime} 假设重新裁剪得到的模型是嵌入的，这一点在模型本身上用计算检查。

\subsection{扇形与锥}\label{sec:fans}

\subsubsection*{设定}
$K$ 是正则的：每个面是由单个线框围成的圆盘，每条边有两个不同的顶点。定理及其证明（附录~\ref{sec:proof51}）只覆盖这种情形。对 B-rep 复形，假设有自然的解读：在接缝边上，E3$'$ 指面的两片；在闭边上，V1 按边端计数；对带内线框的面，P0 要求 $Q_f$ 同胚于该面。但证明不覆盖这些情形。例如在接缝处，同一个面的两片扇形会粘成一个环。这些情形本文没有证明。$\mathrm{dev}^*_e(t)$ 是边界点 $b_{f,e}(t)=\partial P_f\cap\Pi(t)$ 在法平面 $\Pi(t)$ 中相对 $\gamma_e(t)$ 的偏差；它不超过 $\mathrm{dev}_e$。

\subsubsection*{辅助集合}
\begin{itemize}
\item \emph{法圆盘。}对边 $e$，选参数 $0<a_e<b_e<1$ 和 $C^1$ 半径函数 $\rho_e(t)$。$\rho_e$ 在两段短桩 $[t_v,a_e]$、$[b_e,t_w]$ 上以及 $a_e$、$b_e$ 的某个邻域上为常数；这里 $t_v$ 是满足 $\gamma_e(t)\in B(p_v,r_v)$ 的最后一个参数。$D_e(t)$ 是 $\Pi(t)$ 中以 $\gamma_e(t)$ 为圆心、半径 $\rho_e(t)$ 的闭圆盘，$T_e=\bigcup_{t\in[a_e,b_e]}D_e(t)$。
\item \emph{顶点区域。}对顶点 $v$，选 $r_v$。顶点区域 $W_v$ 是球 $B(p_v,r_v)$ 与以下圆盘之并：对 $v$ 处的每条边 $e$，取 $t$ 在 $t_v$ 与 $a_e$（或 $b_e$）之间的圆盘 $D_e(t)$。
\end{itemize}

\subsubsection*{假设}
E1--E3$'$ 要求在 $[t_v,t_w]$ 上成立，包括短桩。
\begin{description}
\item[E1] $\mathrm{dev}^*_e(t)<\rho_e(t)<\min(\tau_e,\operatorname{reach}\gamma_e|_{[a_e,b_e]})$。
\item[E2] 对每个面 $f>e$：
  \begin{enumerate}
  \item $\psi_f<\pi/2$，其中 $\psi_f$ 是管内 $\gamma_e'$ 与 $T P_f$ 的最大夹角；并且 $\rho_e k_f<1$，其中 $k_f=\kappa_f/\cos\psi_f$；
  \item $t\mapsto\pi_e(\beta_{f,e}(t))$ 严格单调。一个充分条件是 $L_e+\kappa_\gamma\mathrm{dev}^*_e<1$，其中 $L_e$ 是 $\beta_{f,e}-\gamma_e$ 关于弧长的 Lipschitz 常数；
  \item \emph{不回返：}$\dist(\gamma_e(t),P_f\setminus C_f)>\rho_e(t)$。这里 $C_f$ 是 $P_f$ 的截弧（从 $b_{f,e}$ 到出口点）之并，其长度至多为 $\ell_f=(\mathrm{dev}^*+\sqrt{\mathrm{dev}^{*2}+\mu\rho^2})/\mu$，$\mu=1-\rho k_f$。
  \end{enumerate}
\item[E3$'$] 两个面截线在 $\partial D_e(t)$ 上的出口点 $k_{f,e}(t)$ 与 $k_{g,e}(t)$ 不同。一个充分条件（E3）是 $\alpha_e(t)>\varepsilon_f+\varepsilon_g$，其中 $\varepsilon=\arcsin\min\{1,(\mathrm{dev}^*_e+k\ell^2/2)/\rho_e\}$，$\alpha_e$ 是两条截线向内切向量的夹角。
\item[V0] $\gamma_e$ 在 $v$ 处的端点以及各面片在 $v$ 处的角点都在 $\inter B(p_v,r_v)$ 内。
\item[V1] (a) $W_v\subseteq\inter\Z(v)$，且 $W_v$ 关于 $p_v$ \emph{严格星形}：对每个 $x\in\partial W_v$，向量 $p_v-x$ 落在 $W_v$ 于 $x$ 处切锥的内部，并且从 $p_v$ 出发的每条射线与 $\partial W_v$ 恰好相交一次。（在连接点处，第二条性质不能由第一条推出，射线可能擦过 $\partial W_v$，所以单独列出。）
(b) \emph{（非局部）}$W_v$ 之外的构造与 $\partial W_v$ 交于一条简单闭曲线 $\lambda_v$。
\begin{itemize}
\item 每条相邻边贡献一个点，每个相邻面贡献一段弧，按 $v$ 的链环的循环顺序排列。
\item 外部各块横截地穿过 $\partial W_v$：沿 $\partial W_v$ 的光滑层（球面和短桩圆柱面），以及沿短桩圆柱面与球面相交的连接曲线。
\item $\lambda_v$ 与这些连接曲线横截相交，并避开 $\partial W_v$ 的角点。
\end{itemize}
\item[S1] (a) 各 $W_v$ 两两不交。(b) 各 $T_e$ 两两不交。(c) 对 $w\notin\partial e$，$\gamma_e$ 和 $T_e$ 不碰 $W_w$。(d) 对 $e\not<f$，$P_f$ 不碰 $T_e$；对 $v\not<f$，$P_f$ 不碰 $W_v$。(e) 在 $\bigcup T\cup\bigcup W$ 之外，各面片块两两不交。(f) 对 $v<e$，$T_e\cap W_v=D_e(a_e)$；$W_v$ 外的每一块与 $W_v$ 只交于 $\lambda_v$。
\item[P0] \emph{（非局部）}每个 $P_f$ 是嵌入的，且 $P_f\cap(\bigcup T\cup\bigcup W)$ 是 $\partial P_f$ 的一个领圈，从而 $Q_f=P_f\setminus\inter(\bigcup T\cup\bigcup W)$ 是一个闭圆盘。
\end{description}
V1(b) 和 P0 不是局部数据的推论。在第~\ref{sec:experiments}~节的模型中，它们只在``带缝隙的立方体''上得到验证，而且是间接的：通过精确检查构造出的曲面是嵌入的。采样检查程序不检验 V1(b)、P0、超出 S1(d) 的 E2.3，也不检验相邻面的 S1(e)。

\begin{theorem}[扇形与锥构造的存在性；带条件]\label{thm:lemma0}
在 E1--E3$'$、V0、V1、S1 和 P0 之下，容差 B-rep 有一个实现。
\end{theorem}
证明在附录~\ref{sec:proof51}。概要：
\begin{itemize}
\item $h$ 沿边跟随 $\gamma_e$，并在法圆盘里放置由半径组成的扇形。
\item 各面片去掉领圈后，与管保持距离。
\item 曲线 $\lambda_v$ 上的锥填满顶点区域。
\item 单射性来自不交的法圆盘、E3$'$、星形性和 S1/P0。
\item 得到的映射是双 Lipschitz 的，所以足够细的割线逼近是一个 PL 实现。
\end{itemize}

\begin{example}
在一个各面被缝隙拉开的立方体上，这个构造已用有理数精确运算完成：216 个三角形，容差区归属和嵌入性都精确检查过（\texttt{compute/lemma0\_box.py}）。
\end{example}

\subsection{重新求交}

\begin{proposition}[平面，三度顶点]\label{prop:B}
设每个面是由单个线框围成的平面圆盘，$P_f\subseteq\Pi_f$，每个顶点恰在三个面上，且这三个平面线性无关。令 $c_v$ 为三个平面的公共点，对边 $e=vw$ 令 $\ell_e=[c_v,c_w]$，令 $\tilde P_f$ 为按 $\partial f$ 的顺序以 $f$ 各顶点的 $c$ 为顶点的多边形。假设：
\begin{itemize}
\item $c_v\in\Z(v)$，$\ell_e\subseteq\Z(e)$，$\tilde P_f\subseteq\Z(f)$；
\item 各 $\tilde P_f$ 是简单多边形，恰好只在公共边和公共顶点处相交。
\end{itemize}
则到 $\bigcup\tilde P_f$ 上的胞腔映射是一个实现。
\end{proposition}
\begin{proof}
它是到 $\bigcup\tilde P_f$ 上的 PL 同胚，且每个胞腔落在自己的容差区里。
\end{proof}

\begin{proposition}[曲面重新求交]\label{prop:Bprime}
设存在一个\emph{重新裁剪的模型}：紧面片 $\tilde P_f$、曲线 $\tilde\gamma_e$ 和点 $c_v$，记 $\Sigma=\bigcup_f\tilde P_f$，满足：
\begin{description}
\item[P] 每个面有 $C^2$ 正则参数化 $X_f\colon U_f\to\R^3$，其中 $U_f$ 是 $\R^2$ 中的开集，或是商 $\R^2/\Lambda_f$（按周期格作商，接缝就这样表示）中的开集。存在紧区域 $D_f\subset U_f$，使 $X_f(D_f)=\tilde P_f$ 且 $X_f$ 在 $D_f$ 上单射。排除奇异参数化，例如球面极点和锥顶。
\item[C] 存在同胚 $\phi\colon|K|\to\Sigma$，把 $K$ 的每个开胞腔映到 $\Sigma$ 的对应开胞腔上。$\Sigma$ 的开胞腔是：点 $c_v$；曲线 $\tilde\gamma_e$ 去掉端点；面片去掉它所有的边曲线（包括接缝）。
\item[R] 令 $G_f\subset D_f$ 为 $f$ 的\emph{边图}：$f$ 所有边曲线（包括接缝；在商参数化中接缝位于 $D_f$ 内部而不在其边界上）在 $X_f$ 下的原像。$G_f$ 是有限图。其顶点是各点 $c_v$ 的原像。其边是简单的 $C^2$ 正则弧，或挂在某个顶点上的环，只在顶点处相交。在 $G_f$ 的每个顶点处，$D_f$ 在相邻两条弧之间的每个扇区的角度都在 $(0,2\pi)$ 内。允许角度为 $\pi$ 的光滑通过，例如光滑边界曲线上的接缝顶点。
\item[B1--B3] $c_v\in\Z(v)$，$\tilde\gamma_e\subseteq\Z(e)$，$\tilde P_f\subseteq\Z(f)$。
\item[T] 在每条边曲线 $\tilde\gamma_e$ 的每个内点处，两个相邻面片（或接缝的两侧）的向内切半平面不同：曲面不会折回自身。在每个顶点 $c_v$ 处，$\Sigma$ 的切锥是一个嵌入的圆盘。该切锥是切扇区 $dX_f(\text{扇区})$ 之并，取遍 $c_v$ 各原像处边图 $G_f$ 的所有扇区。这个条件的意思是：这些扇区只沿公共边界射线相接，并按 $v$ 的链环的循环顺序排列。（结合直到顶点的 $C^1$ 正则性，这也控制了边端点处的边条件。）
\end{description}
则容差 B-rep 有一个实现。
\end{proposition}

\begin{remark}
\begin{itemize}
\item 由 C，$\Sigma$ 是闭曲面的不交并，所以按定义~\ref{def:solid} 它围成一个实体。
\item C 覆盖接缝、闭边和带内线框的面。
\item T 的边部分只要求曲面不折叠。横截相交的面满足它；相切（G$^1$）边也满足，那里两个半平面方向相反。所以只要有一个具备上述性质的重新裁剪模型，这个命题原则上也适用于带圆角的模型。我们的计算是靠延伸曲面再重新求交来构造重新裁剪模型的，这在面相切时会失败（第~\ref{sec:experiments}~节）；这是计算的局限，不是命题的局限。
\end{itemize}
\end{remark}
证明在附录~\ref{sec:proofBprime}。

\begin{remark}[假设的检查]
\begin{itemize}
\item 重新裁剪模型可以通过延伸各面曲面、重建实体来计算（OpenCascade 的 \texttt{BOPAlgo\_MakerVolume}）。然后在采样点上检查假设（\texttt{compute/reintersect\_curved.py}），但假设 T 的顶点部分没有检查。这对圆柱和管很要紧，因为它们的接缝顶点只落在两张曲面上。证明中的 PL 步骤没有数值实现；它是一个存在性论证。
\item 对平面情形，缝合盒子和缝合平板的命题~\ref{prop:B} 的假设也用有理数精确运算复查过，包括嵌入性。
\item 结果见第~\ref{sec:experiments}~节。
\end{itemize}
\end{remark}

\section{一个障碍}\label{sec:obstruction}

\subsection{设定}
\begin{itemize}
\item $f\ne g$ 是两个面，它们的闭胞腔是由单个线框围成的圆盘。
\item 它们的数据是平面的：$P_f\subseteq\Pi_f$，$P_g\subseteq\Pi_g$，两平面不平行。
\item $\lambda=\Pi_f\cap\Pi_g$，$\alpha\in(0,\pi/2]$ 是两平面的夹角。
\item $n_f,n_g$ 是单位法向量，$m_f$ 是 $\Pi_f$ 中与 $\lambda$ 正交的单位向量（$m_g$ 同理）。
\item $\pi_f$ 是到 $\Pi_f$ 的正交投影，$U_f=\bigcup_{e<f}\Z(e)$，它包含 $f$ 各顶点的容差区。
\item $\ell_f$ 是 $f$ 的边界环：按 $\partial f$ 的循环顺序排列的弧 $\beta_{f,e}$，相邻两弧用连接其端点的线段接起来。
\end{itemize}

\begin{theorem}[障碍]\label{thm:obstruction}
设存在 $X\in\lambda$ 和 $\rho>0$，使得对 $J_f=X+[-\rho,\rho]m_f$、$J_g=X+[-\rho,\rho]m_g$ 有：
\begin{description}
\item[O0] 对 $f$ 的每条边 $e$，$\beta_{f,e}\subseteq\Z(e)$；在 $f$ 的每个顶点 $v$ 处，$f$ 在 $v$ 处的两段弧的端点都在 $\Z(v)$ 中；对 $g$ 同样。
\item[O1] $J_f\cap\ell_f=\emptyset$，且 $\ell_f$ 在 $\Pi_f$ 中绕 $J_f$ 上各点的环绕数为奇数；对 $g$ 同样。
\item[O2] $J_f\cap\cl\,\pi_f(U_f)=\emptyset$，$J_g\cap\cl\,\pi_g(U_g)=\emptyset$。
\item[O3] 对 $f$ 或 $g$ 的每条边 $e$，$\Z(e)$ 单连通。
\item[O4] $\rho\sin\alpha>\tau_f\cos\alpha+\tau_g$ 且 $\rho\sin\alpha>\tau_g\cos\alpha+\tau_f$。
\end{description}
则该容差 B-rep 没有实现。
\end{theorem}

\begin{proof}
设 $h$ 是一个实现，令 $\varphi=\pi_f\circ h|_{\bar f}$。

\emph{第 1 步：在 $J_f$ 上方，$h(\bar f)$ 落在薄板内。}
\begin{itemize}
\item 设 $y\in J_f$。有 $h(\bar f)\subseteq\Z(f)=N_{\tau_f}(P_f)\cup U_f$，由 O2，$U_f$ 中没有点在 $y$ 的正上方（投影为 $y$）。
\item 由于 $P_f\subseteq\Pi_f$，$\pi_f^{-1}(y)\cap N_{\tau_f}(P_f)\subseteq\sigma_y:=y+(-\tau_f,\tau_f)n_f$。
\item 所以 $h(\bar f)\cap\pi_f^{-1}(y)=h(\bar f)\cap\sigma_y$。这些点是 $f$ 的内点的像，因为 $h(\partial f)\subseteq U_f$。
\end{itemize}

\emph{第 2 步：度为一。}
\begin{itemize}
\item 在 $\partial f$ 上参数化 $\ell_f$：每个顶点周围的一小段弧 $A_v$ 映到连接线段，每条边的其余部分映到 $\beta_{f,e}$。把 $A_v$ 缩小，直到 $h(A_v)\subseteq\Z(v)$。
\item 在 $U_f$ 中构造从 $h|_{\partial f}$ 到这个参数化的同伦：
  \begin{itemize}
  \item 在 $A_v$ 上，用凸球 $\Z(v)$ 内的直线同伦；
  \item 在边 $e$ 的其余部分上，正方形的边界映入 $\Z(e)$（由 O0、单调性和 $h(\bar e)\subseteq\Z(e)$），由 O3 可延拓到整个正方形上。
  \end{itemize}
\item 由 O2，投影后的同伦避开每个 $y\in J_f$。所以 $\varphi|_{\partial f}$ 与 $\ell_f$ 绕 $y$ 的环绕数相同，由 O1 为奇数。因此 $\varphi$ 在 $y$ 处的模 2 度为 $1$。
\item 取 $\bar f$ 的一个使 $h$ 线性的三角剖分。对 $y\notin\varphi(\text{1-骨架})$（一旦第 3 步中的截面取成一般位置，这只排除 $J_f$ 上的有限个点），$\#(h(\bar f)\cap\sigma_y)$ 为奇数。
\end{itemize}

\emph{第 3 步：一条横穿的弧。}
\begin{itemize}
\item 令 $\Sigma$ 为过 $X$ 且与 $\lambda$ 正交的平面；它包含 $n_f,m_f,n_g,m_g$。
\item 沿 $\lambda$ 移动 $X$，把 $\rho$ 略为缩小，取一般位置的值；O1、O2、O4 都是开条件。一般位置指：$\Sigma$ 不含三角剖分的顶点、与其边横截相交，且直线 $\langle x-X,m_f\rangle=\pm\rho$、$\langle x-X,m_g\rangle=\pm\rho$ 不含顶点。
\item 在 $\Sigma$ 中，矩形 $R_f=J_f+[-\tau_f,\tau_f]n_f$ 与 $h(\bar f)$ 交于一个紧 1-流形 $M_f$。
\item $M_f$ 与 $h(\partial f)$ 不交，因为 $\pi_f(R_f)=J_f$ 且 O2 成立。由第 1 步，它与 $R_f$ 的长边不交。所以它的边界落在两条短边上。
\item 一般的纤维 $\sigma_y$ 与 $M_f$ 相交奇数次（第 2 步）。圆周分支，以及两端都在同一条短边上的弧，与它相交偶数次。
\item 所以 $M_f$ 中有一条弧 $a_f$ 连接 $R_f$ 的两条短边。对 $g$ 同样有弧 $a_g$。
\end{itemize}

\emph{第 4 步：两条弧相交。}
\begin{itemize}
\item 在 $\Sigma$ 中用坐标 $u=\langle x-X,n_f\rangle$、$w=\langle x-X,n_g\rangle$。
\item 在 $R_f$ 的短边上，由 O4，$|w|\ge\rho\sin\alpha-\tau_f\cos\alpha>\tau_g$，且两条短边上符号相反。所以 $a_f$ 在平行四边形 $Q=\{|u|\le\tau_f,|w|\le\tau_g\}$ 中含有一段连接 $w=-\tau_g$ 与 $w=\tau_g$ 的子弧。对称地，$a_g$ 在 $Q$ 中含有一段连接 $u=-\tau_f$ 与 $u=\tau_f$ 的子弧。
\item 由 O4，$Q$ 落在两个矩形中。
\item 正方形中分别连接两组对边的两条道路必相交（Brouwer 不动点定理的标准推论；未单独对照原文）。所以 $h(\bar f)\cap h(\bar g)\ne\emptyset$，交点同时是 $f$ 和 $g$ 的内点的像（第 1 步）。这与单射性矛盾。
\end{itemize}
\end{proof}

\begin{remark}
\begin{itemize}
\item \emph{哪些假设起约束作用。}面容差很小（OpenCascade 默认 $10^{-7}$），所以 O4 很宽松。真正起约束作用的是 O1 和 O2：两张面片必须在边区和顶点区的``影子''之外相交。
\item \emph{曲面面片。}证明在第 1、2 步用到平面性。对 $\tau_f$ 小于其 reach 的 $C^2$ 面片，用管状坐标应能得到相应的版本；本文没有写出。
\item \emph{用到而未单独对照原文的标准事实：}PL 映射的模 2 度及其与环绕数的关系、PL 映射与一般平面的一般位置、线段与 PL 圆周交点的奇偶性、环绕数的同伦不变性、神经引理（用于计算中的 O3 检验），以及上面引用的正方形引理。
\end{itemize}
\end{remark}

\section{OpenCascade 实验}\label{sec:experiments}

\subsection{设置}
\begin{itemize}
\item \emph{软件。}所有实验都通过 \texttt{cadquery-ocp} 使用 OpenCascade 8.0.1。脚本和结果在代码仓库中~\cite{repo}。
\item \emph{两项检查。}``BRepCheck''指 \texttt{BRepCheck\_Analyzer(shape).IsValid()}。``BOP 检查''指 \texttt{BRepAlgoAPI\_Check(shape, true, true)}，它还检查几何和自相交。
\item \emph{缝合模型}模拟导入的数据。把一个精确零件的每个面随机平移一个长度在 $g/2$ 到 $g$ 之间的向量，再用 \texttt{BRepBuilderAPI\_Sewing} 以容差 $4g$ 缝合回去。缝合把每条边的容差设为它自己偏差估计值的 $1.05$ 倍，即我们实测偏差的 $1.048$--$1.05$ 倍。
\item \emph{范围。}这只是一个内核版本里的几个零件，不是对工业模型的普查。
\end{itemize}

\subsection{检查判定不了存在性和唯一性}
表~\ref{tab:occ} 是汇总。

\begin{table}[h]
\centering\small
\begin{tabular}{>{\raggedright}p{4.4cm}cc>{\raggedright\arraybackslash}p{5.0cm}}
\toprule
模型 & BRepCheck & BOP 检查 & 实际情况 \\
\midrule
扭转的三棱柱，两个面的边界相互链环（$\tau=0.081$） & 有效 & 自相交 & 无实现（链环数 $-1$） \\
厚壁管，容差为孔半径的 $1.3$ 倍 & 有效 & 有效 & 两个实现，实体不同胚（结行列式 $3$ 对 $1$） \\
缝合楔形体，$30^\circ$ 和 $10^\circ$，$g=0.01$ & 有效 & 有效 & 无实现（定理~\ref{thm:obstruction}） \\
缝合平板，厚 $0.05$，$g=0.01$ & 有效 & 无效 & 可实现（命题~\ref{prop:B}） \\
带圆角的 L 形，$\tau=0.1$ & 有效 & 自相交 & $\tau+\mathrm{defl}<M_{PL}$ 的认证 $\wfs$ 下界 \\
长圆柱，$\tau=3.3$ & 有效 & 有效 & $\tau\ge\wfs$（名义）（给出了临界点） \\
\bottomrule
\end{tabular}
\caption{OpenCascade 的检查与理论结论的对照。}\label{tab:occ}
\end{table}

\subsubsection*{扭转的三棱柱}
\begin{itemize}
\item 两个不相邻的三角形面 $A,B$ 的边界环的链环数为 $-1$（由交叉计数得到；数值 Gauss 积分结果一致）。两个环相距 $0.4037>2\tau=0.1615$，所以它们的边区与顶点区之并 $T_A$、$T_B$ 不交。
\item 在任一实现中，每个 $h(e)$ 落在凸胶囊 $C_e=N_\tau(e)$ 内，它包含两端的球（这里所有容差相等）。把 $h(e)$ 在 $C_e$ 内保持端点不动地同伦到一条线段，把顶点在球内同伦到名义顶点。这就在 $T_A$ 内、因而在 $T_B$ 的补集内，给出 $h(\partial A)$ 到 $\partial A$ 的同伦；$B$ 同样。
\item 链环数在这样的同伦下不变，所以 $\mathrm{lk}(h(\partial A),h(\partial B))=-1$，$h(\partial A)$ 必与圆盘 $h(\bar B)$ 相交。这与嵌入性矛盾。
\end{itemize}

\subsubsection*{打结的管}
\begin{itemize}
\item 一个实现保持名义几何。另一个把内表面换成一根半径 $0.005$ 的细管，其中心线沿一条拉长的三叶结走。
\item 脚本检查了容差区（最大距离 $0.983\le1.3$）、嵌入性（$\varepsilon\kappa_{\max}=0.31$；在窗口 $0.04<\pi/\kappa_{\max}=0.0504$ 之外的非局部距离 $0.036>2\varepsilon$）和结（行列式 $3$）。
\item 第二个实体是球中一条打结弧的补空间，其基本群不交换：行列式为 $3$ 的结有非平凡的 Fox 3-染色，而图上的非平凡 3-染色通过 Wirtinger 表示给出结群到不交换群 $S_3$ 的满射（标准事实；未单独对照原文）。名义实体是实心环。所以两个实体不同胚。
\item 这说明定理~\ref{thm:uniqueness} 的领圈假设不能去掉。这里 $\wfs\le1<1.3$。
\end{itemize}

\subsubsection*{缝合楔形体}
\begin{itemize}
\item 在缝合楔形体的刃边处，两张平面面片沿一条直线相交，这条线落在边管之外，超出量为 $\tau_e$ 的 $49\%$（$30^\circ$）和约 $3$ 倍（$10^\circ$）。
\item 定理~\ref{thm:obstruction} 的假设成立。对 $30^\circ$ 楔形体，在所搜索的位置上可取的最大 $\rho$ 是 $1.6\cdot10^{-3}$，而 O4 只需要 $3.7\cdot10^{-7}$；对 $10^\circ$ 楔形体是 $1.6\cdot10^{-2}$。
\item 一次独立的精确有理数运算检查，确认两个楔形体都满足 O0--O4。
\item 在缝合盒子、缝合平板以及精确盒子、精确楔形体上，检查程序没有找到相交见证，这与它们可实现一致。
\end{itemize}

\subsection{哪种构造适用}

\begin{table}[h]
\centering\small
\begin{tabular}{>{\raggedright}p{4.0cm}>{\raggedright}p{2.8cm}>{\raggedright}p{2.6cm}>{\raggedright\arraybackslash}p{4.0cm}}
\toprule
模型 & 扇形（定理~\ref{thm:lemma0}） & 重新求交 & 重新求交的余量（顶点 / 边 / 面） \\
\midrule
精确盒子、圆柱、管、槽、带圆角盒子、$10^\circ$ 楔形体 & 通过 & 适用处通过 & --- \\
L 形，所有边倒圆角 0.5 & 用变半径通过$^*$ & 不适用 & --- \\
融合的 L 形，内圆角 0.2 & 失败（2 个尖点顶点，9 处 V1(a) 失败）；合并共面面后通过 & 不适用 & --- \\
缝合盒子 $g=0.01$ & 在顶点处失败 & 通过 & 45\% / 13\% / 4.9\% \\
缝合平板 $g=0.01$ & 在顶点处失败 & 通过 & 45\% / 13\% / 5.9\% \\
缝合圆柱 $g=0.01$ & 在顶点处失败 & 通过（曲面） & 55\% / 35\% / 5.0\% \\
缝合圆柱 $g=0.05$ & 在顶点处失败 & 通过（曲面） & 55\% / 11\% / 8.9\% \\
缝合管 $g=0.01$ & 在顶点处失败 & 通过（曲面） & 53\% / 16\% / 5.1\% \\
缝合楔形体 & 在顶点处失败 & 失败（边落在管外） & 无实现（定理~\ref{thm:obstruction}） \\
缝合带圆角盒子 $g=0.001$ & 在顶点处失败 & 计算不适用（相切面） & 未解决 \\
\bottomrule
\end{tabular}
\caption{在精确模型和缝合模型上的构造（采样检查，双精度）。$^*$暂定：检查程序还没有实现顶点附近的常半径。面余量取采样检查程序的值与一次用精确投影的独立密集复查（\texttt{results/reintersect\_B3\_dense.txt}）中较小者；缝合平板只有密集复查的值。}\label{tab:constr}
\end{table}

\begin{itemize}
\item \emph{扇形为什么在缝合数据上失败。}扇形构造在缝合模型的顶点处失败。缝合把 $\tau_e$ 设为 $1.05\,\mathrm{dev}_e$，而管半径必须超过偏差；但在盒子的角上，三根半径都大于 dev 的管，没法互不相交地放进一个半径约为 $\tau_e$ 的顶点球里。
\item \emph{重新求交。}重新求交避开了这个问题。我们计算重新裁剪模型的方法是把延伸后的曲面重新求交，这需要面横截相交。``通过（曲面）''只指假设 T 的边部分；顶点部分没有检查。
\item \emph{带相切边的缝合模型仍未解决。}在缝合带圆角盒子上，OpenCascade 把每个顶点容差设为相邻边容差的最大值，面片角点离顶点最远达 $0.85\,\tau_v$，我们的顶点构造都放不下。
\end{itemize}

\subsection{其他观察}
\begin{itemize}
\item \emph{厚度经验估计。}不相邻面间最小距离的一半，是常用的安全容差估计。它高估了圆柱（5 对 $\wfs=3$），对一个内圆角为 $0.05$ 的 L 形则低估了约 $97$ 倍。
\item \emph{偏置。}在默认模式下，\texttt{BRepOffsetAPI\_MakeOffsetShape} 对一个带圆角盒子、向内偏置距离大于圆角半径但仍在 $\wfs$ 安全范围内时，返回了体积错误的实体（误差从 $44\%$ 到上千倍，每一个都是六个壳而不是一个），而 BRepCheck 都接受了。OpenCascade~8.0.1 的源码里有一个可能的机理：圆柱面、球面或圆环面形式的圆角面，偏置半径变成负值时，其偏置曲面被标为 \texttt{BRepOffset\_Reversed}，而 \texttt{BRepOffset\_MakeOffset::BuildOffsetByArc} 跳过这类面（源码注释：``Degenerated or returned faces are not stored''）。算法的其余部分我们没有追查。这类失败看来与容差无关，像是``逐面偏置''在面必须消失时的失败。
\item \emph{BRepCheck 为什么不管容差大小。}在 OpenCascade~8.0.1 中，\texttt{BRepCheck\_Analyzer} 的文档注释把状态 \texttt{BRepCheck\_InvalidToleranceValue} 对顶点、边、面都标为``NYI''（尚未实现）。在 \texttt{BRepCheck\_Analyzer.cxx} 里，本该给面设置这个状态的标志初始化为假后从未赋值，其他文件也不设置这个状态。所以 BRepCheck 从不报告容差过大。（这一点最先在作者的配套著作项目中发现；源码我们自己核对过。）
\item \emph{公开文档。}我们在 OpenCascade 的公开文档以及 ACIS、Parasolid 的第三方文档中，都没有找到``检查容差不改变拓扑''的功能。这只说明文档里没有提到。
\end{itemize}

\section{局限与未解决问题}\label{sec:open}

\begin{enumerate}
\item \textbf{证明背后的引用。}证明中引用的文献结果都已对照原文核对。剩下的是每处都注明``未单独对照原文''的标准事实。
\item \textbf{扇形构造中的 B-rep 复形。}定理~\ref{thm:lemma0} 只对正则复形证明；接缝、闭边和带内线框的面需要另外论证。（命题~\ref{prop:Bprime} 覆盖了它们。）
\item \textbf{球面。}定理~\ref{thm:uniqueness} 不含亏格为 0 的壳，Waldhausen 的水平化论证对它们不适用。PL Schoenflies 定理~\cite[Thm.~17.12]{Moise77} 是处理球面情形的自然工具；论证没有写出。
\item \textbf{嵌入。}我们证明的是实现的像环境同痕，而不是嵌入映射作为映射同痕。两者相差曲面的一个自同胚。
\item \textbf{紧顶点。}带相切边、使用 OpenCascade 顶点容差（等于相邻边容差的最大值）的缝合模型是否有实现，我们不知道。两个方向都未解决：一种同时利用边管和球的顶点构造，或者障碍定理的顶点版本。
\item \textbf{相切边的重新裁剪模型。}命题~\ref{prop:Bprime} 允许相切（G$^1$）边和混合顶点，但我们计算重新裁剪模型的方法是把延伸后的曲面重新求交，在面相切时会失败。如何为这类数据构造重新裁剪模型，仍未解决。
\item \textbf{曲面版障碍。}定理~\ref{thm:obstruction} 只对平面面片陈述。
\item \textbf{计算的严格性。}
  \begin{itemize}
  \item $\wfs$ 下界和假设检查用的是双精度和采样。例外：带缝隙立方体的构造，以及缝合楔形体、缝合盒子、缝合平板的精确复查。
  \item 对工业模型，需要区间算术和编译型实现。
  \end{itemize}
\item \textbf{形式化。}一般情形的承载引理（可缩承载集）、Waldhausen 定理和各存在性构造都没有形式化。
\end{enumerate}

\appendix
\renewcommand{\proofname}{Proof}
\section{A curved-triangle lemma}\label{sec:curvedtri}

Appendices~\ref{sec:proof51} and~\ref{sec:proofBprime} need maps of triangles with one curved side that reproduce a
\emph{prescribed} parametrization of that side. This is what makes maps built separately on two adjacent faces agree
on their common edge. We give an explicit construction.

\begin{lemma}\label{lem:curvedtri}
Let $\hat c\colon[t_0,t_1]\to\R^2$ be a $C^2$ regular curve, $a_1=\hat c(t_0)$, $a_2=\hat c(t_1)$, and $a_3$ a point.
\begin{itemize}
\item Put $\psi(s)=\hat c(t_0+s(t_1-t_0))-((1-s)a_1+sa_2)$ for $s\in[0,1]$.
\item In barycentric coordinates $(\lambda_1,\lambda_2,\lambda_3)$ of the reference triangle $\Delta$, put
$\mu=\lambda_1+\lambda_2$ and
\[
F(\lambda)=\lambda_1a_1+\lambda_2a_2+\lambda_3a_3+\mu^2\,\psi(\lambda_2/\mu)\qquad(F=a_3\text{ at }\mu=0).
\]
\end{itemize}
Then:
\begin{enumerate}
\item[(a)] on the side $\lambda_3=0$, $F$ is $\lambda_2\mapsto\hat c(t_0+\lambda_2(t_1-t_0))$;
\item[(b)] on the sides $\lambda_1=0$ and $\lambda_2=0$, $F$ is affine;
\item[(c)] $F$ is $C^1$ on the closed triangle $\Delta$, including at its corners.
\end{enumerate}
Moreover, consider a family of such triangles in which the chord triangles $(a_1,a_2,a_3)$ are shape-regular with
diameter $h\to0$, and $|t_1-t_0|$ times the speed of $\hat c$ is $O(h)$, with uniform $C^2$ bounds on $\hat c$. Then for
small $h$:
\begin{enumerate}
\item[(d)] $F$ is a $C^1$ diffeomorphism of $\Delta$ onto the region bounded by the arc $\hat c([t_0,t_1])$ and the
segments $[a_2,a_3]$ and $[a_3,a_1]$, and $DF=A(I+O(h))$, where $A$ is the derivative of the affine part.
\end{enumerate}
\end{lemma}

\begin{proof}
(a) On $\lambda_3=0$ we have $\mu=1$, so $F=(1-\lambda_2)a_1+\lambda_2a_2+\psi(\lambda_2)=\hat c(t_0+\lambda_2(t_1-t_0))$.

(b) $\psi(0)=\psi(1)=0$. On $\lambda_2=0$ the correction is $\mu^2\psi(0)=0$. On $\lambda_1=0$ it is
$\mu^2\psi(1)=0$ (for $\mu>0$).

(c) For $\mu>0$, $F$ is $C^1$ (indeed $C^2$). Write $G=\mu^2\psi(s)$ with $s=\lambda_2/\mu$. Then
$dG=2\mu\psi(s)\,d\mu+\mu^2\psi'(s)\,ds$ with $|ds|=O(1/\mu)$, so $|dG|=O(\mu)\to0$ as $\mu\to0$. Hence $F$ is $C^1$ at
$a_3$ as well, with $DF(a_3)=A$. At the other two corners $\mu=1$, and $\psi$ is $C^2$ on $[0,1]$.

(d) \emph{Size of the correction.} $\psi''(s)=(t_1-t_0)^2\hat c''$ is $O(h^2)$, and $\psi(0)=\psi(1)=0$, so $|\psi|$ and
$|\psi'|$ are $O(h^2)$. Hence $|DG|=O(h^2)$ on $\Delta$, in coordinates of the unit reference triangle.

\emph{Injectivity.} By shape regularity $\|A^{-1}\|=O(1/h)$, so $DF=A+DG=A(I+O(h))$. For $x,y\in\Delta$ (which is
convex),
\[
|F(x)-F(y)|\ge|A(x-y)|-Ch^2|x-y|\ge(ch-Ch^2)|x-y|>0
\]
for small $h$. So $F$ is injective, and its differential is invertible.

\emph{Image.} $F$ maps $\partial\Delta$ homeomorphically onto the closed curve formed by the arc and the two segments.
That curve is simple for small $h$: the arc stays within $O(h^2)$ of its chord, and the two segments make angles with
the chord that are bounded below by shape regularity. An injective continuous map of a closed disk whose boundary goes
onto a Jordan curve has as image the closed region bounded by that curve.
\end{proof}

\emph{Remark.} The classical blending map with $\mu\,\psi(\lambda_2/\mu)$ in place of $\mu^2\psi(\lambda_2/\mu)$ is
exact on the curved side and affine on the straight sides. But it is only Lipschitz at $a_3$: its derivative there
depends on the direction of approach. The factor $\mu^2$ is what gives (c).

\emph{Standard facts used without separate checking against a source:} the Jordan curve theorem, and that an injective
continuous map of a closed disk onto a set bounded by its image boundary curve has the closed Jordan region as its
image (invariance of domain).

\section{Proof of Theorem~\ref{thm:lemma0}}\label{sec:proof51}

We use the notation and hypotheses of Section~\ref{sec:fans}. The \emph{exit point} $k_{f,e}(t)$ is the point where the
section of $P_f$ in $\Pi(t)$, followed from $b_{f,e}(t)$, leaves $D_e(t)$; E2 makes it unique. $m>0$ denotes the
minimum zone margin: $\rho_e<\tau_e$ and $W_v\subseteq\inter\Z(v)$, and by compactness all pieces built below lie at
distance at least $m$ inside the zones of their cells.

\subsection{Construction of $h$}
\begin{enumerate}
\item $h(v)=p_v$.
\item On $[a_e,b_e]$, $h(e)=\gamma_e$. In each $D_e(t)$, the radius $[\gamma_e(t),k_{f,e}(t)]$ is assigned to $f$. The fan
strip is $F_{f,e}=\bigcup_{t\in[a_e,b_e]}[\gamma_e(t),k_{f,e}(t)]$.
\item Outside the vertex regions, $h(f)=Q_f\cup\bigcup_{e<f}F_{f,e}$, where $Q_f=P_f\setminus\inter(\bigcup T\cup\bigcup W)$.
\item Inside $W_v$, take the cone from $p_v$ over $\lambda_v$: each $x\in\lambda_v$ is joined to $p_v$ by a segment. The
cone over the point of $\lambda_v$ contributed by $e$ is $h(e)$ between $v$ and $a_e$; the cone over the arc contributed
by $f$ is a sector of $h(f)$.
\end{enumerate}

\subsection{$h$ is a continuous injective cellular map into the zones}
\emph{Pieces and their supports.} $Q_f\subseteq P_f$ lies outside $\inter(\bigcup T\cup\bigcup W)$. $F_{f,e}\subseteq T_e$.
The cone sectors and edge stubs lie in $W_v$.

\emph{Injectivity inside $T_e$.}
\begin{itemize}
\item The disks $D_e(t)$ lie in the normal planes and have radius below the reach, so they are pairwise
disjoint~\cite{Federer59}; this does not depend on $\rho_e$ being constant.
\item Inside one disk, the construction consists of the centre and two radii to the distinct points $k_{f,e}(t)$,
$k_{g,e}(t)$ (E3$'$). They meet only at the centre.
\item No other piece enters $T_e$. Non-incident pieces are excluded by S1(b)--(d) and S1(f). The patches of $f>e$
enter $T_e$ only in the collar of $\partial P_f$ (P0 and E2.3), and that collar is removed.
\end{itemize}

\emph{Injectivity inside $W_v$.}
\begin{itemize}
\item The cone map $(s,x)\mapsto p_v+s(x-p_v)$, for $s\in(0,1]$ and $x\in\lambda_v$, is injective. Each ray from $p_v$
meets $\partial W_v$ once (V1(a)), and $\lambda_v$ is a simple curve on $\partial W_v$ (V1(b)).
\item No other piece enters $W_v$: S1(a), (c), (d) for non-incident pieces, and S1(f) for incident ones.
\end{itemize}

\emph{Injectivity elsewhere.} Outside $\bigcup T\cup\bigcup W$ the pieces are the $Q_f$. They are embedded (P0) and
pairwise disjoint (S1(e)).

\emph{Continuity and cell structure.}
\begin{itemize}
\item \emph{The pieces fit.} $Q_f$ and $F_{f,e}$ agree along $k_{f,e}$. The fans and $Q_f$ agree with the cone
along $\lambda_v$, which is by V1(b) exactly their trace on $\partial W_v$.
\item \emph{Order.}
  \begin{itemize}
  \item By P0, $\partial Q_f$ is one circle made of the curves $k_{f,e}$ and arcs $A_{f,v}$ on the spheres $\partial W_v$.
  \item By V1(b) there is exactly one arc per vertex of $f$, joining $k_{f,e}(a_e)$ to $k_{f,e'}(a_{e'})$.
  \item So these pieces form a 2-regular graph isomorphic to $\partial f$, whose only cycle is $\partial Q_f$; the
  cyclic order is forced.
  \end{itemize}
\item \emph{Exact intersections.}
  \begin{itemize}
  \item $Q_f\cap F_{f,e}=k_{f,e}$, by E2.3 and S1(f).
  \item $(Q_f\cup\text{fans})\cap\text{sector}_{f,v}=A_{f,v}$, by star-shapedness and S1(f).
  \item Gluing two disks along a boundary arc that is their whole intersection gives a disk. By induction, $h(f)$ is
  a closed disk with $\partial h(f)=\bigcup_{e<f}h(e)$.
  \end{itemize}
\item \emph{$h$ as a map.} Fix a homeomorphism $\partial f\to\partial h(f)$ that respects the cells, and extend it
over $f$ through a chart of the disk $h(f)$ (Alexander's trick).
  \begin{itemize}
  \item On the part away from the vertex regions, take the extension piecewise $C^1$, by straightening the $C^1$
  boundary of $Q_f$ in the parameter domain of $P_f$.
  \item On each cone sector, take it radial: the sector of $\bar f$ at $v$ is the cone over an arc of the link, and a
  point at radius $s$ over a link point $y$ goes to $p_v+s(\phi(y)-p_v)$. Here $\phi$ is a $C^1$ parametrization of
  the arc $A_{f,v}$ of $\lambda_v$.
  \end{itemize}
\end{itemize}

\emph{Zones.}
\begin{itemize}
\item $Q_f\subseteq P_f\subseteq\Z(f)$.
\item $F_{f,e}\subseteq T_e\subseteq N_{\rho}(\gamma_e)\subseteq\Z(e)\subseteq\Z(f)$.
\item The cone parts lie in $W_v\subseteq\Z(v)$, which is contained in $\Z(e)$ and $\Z(f)$ for incident $e$ and $f$.
\item $h(\inter e)\subseteq\gamma_e\cup W_v\subseteq\Z(e)$.
\item All inclusions hold with margin at least $m$.
\end{itemize}

\subsection{The PL step}
\emph{$h$ is bi-Lipschitz.} $h$ is injective on a compact set, so it suffices that it is locally bi-Lipschitz.
\begin{itemize}
\item \emph{At the creases,} the opening angle is positive.
  \begin{itemize}
  \item Along $k$, between fan and patch: the fan leaves $k$ in direction $-\hat r$, and the patch section leaves with
  $c'\cdot\hat r>0$ (E2), so the angle exceeds $90^\circ$.
  \item Along $\gamma_e$, between two fans: E3$'$ and compactness.
  \item Along $\partial W_v$, between the cone and the outside pieces: strict star-shapedness.
  \end{itemize}
\item \emph{Where three or more pieces meet} ($k_{f,e}(a_e)$, $\gamma_e(a_e)$, corners of $\partial W_v$), the tangent
cone is an embedded PL disk.
  \begin{itemize}
  \item At $k(a_e)$: the fan and $Q_f$ tangent sectors lie in the closed quarter-space
  $\{\langle\cdot,\gamma'\rangle\le0,\langle\cdot,\hat r\rangle\le0\}$. This uses $\rho_e'(a_e)=0$: $\rho_e$ is constant
  near $a_e$. The cone sectors lie inside it, with $-(x-p_v)$ in its interior. The only shared rays are $-\hat r$ and
  the side-wall tangent, each between adjacent sectors.
  \item At $\gamma(a_e)$: the fans lie in $\{\langle\cdot,\gamma'\rangle\ge0\}$ and the sectors in
  $\{\langle\cdot,\gamma'\rangle\le0\}$. They share only the rays $\hat r_f\ne\hat r_g$ (E3$'$).
  \end{itemize}
\item \emph{At the apex $p_v$,} $h$ is homogeneous of degree one on each sector. It is bi-Lipschitz there because
$\lambda_v$ is transverse to the rays from $p_v$ and $|x-p_v|$ is bounded below on $\lambda_v$.
\end{itemize}

\emph{Secant approximation.} Take a shape-regular triangulation of $|K|$ with these properties:
\begin{itemize}
\item it refines the decomposition into pieces;
\item it has the vertices of $K$ as vertices;
\item on each cone sector it is radial: its triangles there have the apex as a vertex and a chord of the link as the
opposite side;
\item $h$ is $C^1$ on each closed triangle away from the apexes.
\end{itemize}
Let $S$ be the secant map, linear on each triangle and equal to $h$ at the vertices.
\begin{itemize}
\item \emph{Away from the apexes.} $h$ is $C^1$ on each closed triangle. By the standard estimate,
$\operatorname{Lip}(S-h)\le C\,\omega(Dh,\mathrm{mesh})/\sin\theta_{\min}$, where $\omega$ is a modulus of continuity and
$\theta_{\min}$ the smallest angle. This tends to $0$ with the mesh size.
\item \emph{On a radial triangle at the apex,} with base chord $[y_0,y_1]$: write $\hat\phi$ for the linear interpolation
of $\phi$ on the chord. Then
\[
(S-h)(p_v+s(y-p_v))=s\,(\hat\phi(y)-\phi(y)),
\]
which is homogeneous of degree one.
  \begin{itemize}
  \item Its radial derivative is $\hat\phi(y)-\phi(y)$, of size at most $\omega(\phi',\mathrm{mesh})\cdot\mathrm{mesh}$.
  \item Its derivative along the chord, divided by the distance $s|y-p_v|$ to the apex, is at most
  $|\hat\phi'-\phi'|/|y-p_v|\le\omega(\phi',\mathrm{mesh})/\min|y-p_v|$.
  \item So $\operatorname{Lip}(S-h)\to0$ uniformly on the cone sectors.
  \end{itemize}
  This uses only that $\phi$ is $C^1$. It closes the gap noted in earlier versions, where the estimate was stated
  with a curvature of $\lambda_v$, which is only $C^1$.
\item \emph{Conclusion.} Since $\operatorname{Lip}(S-h)\to0$ and $|S-h|\to0$, for a fine enough mesh:
  \begin{itemize}
  \item $\operatorname{Lip}(S-h)$ is below the bi-Lipschitz constant of $h$, so $S$ is injective;
  \item $|S-h|<m$, so $S(\inter\sigma)\subseteq\Z(\sigma)$.
  \end{itemize}
  $S$ is PL, so it is a realization. $\square$
\end{itemize}

\emph{Standard facts used without separate checking against a source:} disjointness of normal disks below the
reach~\cite{Federer59}, the implicit function theorem (for the $C^1$ dependence of $k_{f,e}$ on $t$), Alexander's
trick, gluing disks along boundary arcs, and the secant-approximation estimate for $C^1$ maps on shape-regular
triangulations.

\section{Proof of Proposition~\ref{prop:Bprime}}\label{sec:proofBprime}

The idea: take the re-trimmed surface $\Sigma$ itself as the model of $|K|$ and the identity as a topological
realization, then replace it by a fine secant approximation. The steps parallel the PL step of
Appendix~\ref{sec:proof51}, but there are no fans or cones.

\subsection{Margin}
By C, the identity $\mathrm{id}\colon\Sigma\to\Sigma$, together with the cell isomorphism, maps each closed cell of $K$
onto the corresponding closed cell of $\Sigma$. By B1--B3, every closed cell lies in the zone of its cell. Closed cells
are compact and zones are open, so there is $m>0$ such that $N_m$ of every closed cell lies in its zone.

\subsection{A curved triangulation of $\Sigma$}
\begin{itemize}
\item \emph{Edges.} On each curve $\tilde\gamma_e$ fix one $C^2$ regular parametrization (say by arc length). Choose
vertices on it, including its end points.
\item \emph{Faces.} In the parameter domain of each $S_f$, the region of $\tilde P_f$ is bounded by the preimages of the
$\tilde\gamma_e$; they carry the pulled-back parametrizations and vertices. Across a seam the parameter domain is
glued periodically, so both sides of a seam use the same vertices.
\item \emph{Curved triangles.} Subdivide each region into curved triangles with at most one curved edge each. The curved
edges are pieces of the boundary curves, with their pulled-back parametrizations; the corners of the region (the
preimages of the $c_v$) are vertices; and the triangles are shape-regular in the parameter domain. This is possible for
a fine enough subdivision because the boundary curves are $C^2$ and meet at positive angles (R).
\item \emph{Maps.} Map each straight reference triangle onto its curved triangle by a transformation that reproduces the
prescribed parametrization on the curved edge and is affine on the straight edges~\cite{Lenoir86}
\unverified{regularity up to the corners and invertibility of Lenoir's transformation}. Compose with the
parametrization of $S_f$.
\end{itemize}
This gives an abstract simplicial complex $T$ and a homeomorphism $h_0\colon|T|\to\Sigma$.
\begin{itemize}
\item $T$ subdivides the cell structure of $\Sigma$, so by C it is a regular subdivision of $K$. It gives $|K|$ its PL
structure (Section~\ref{sec:setting}).
\item On each closed triangle, $h_0$ is $C^1$ up to the boundary, with injective differential.
\item On a shared edge, the two adjacent triangle maps agree:
  \begin{itemize}
  \item an interior straight edge of one face is affine in the parameter domain on both sides;
  \item an edge on $\tilde\gamma_e$ reproduces the same 3D parametrization from both faces (or from both sides of a seam).
  \end{itemize}
\end{itemize}

\subsection{$h_0$ is bi-Lipschitz}
Use the intrinsic length metric of $|T|$. $h_0$ is injective by C, and $|T|$ is compact, so it suffices that $h_0$ is
locally bi-Lipschitz.
\begin{itemize}
\item \emph{Inside a face,} including across interior edges of the subdivision: $h_0$ is a composition of a
bi-Lipschitz map onto the parameter region with the regular parametrization of $S_f$.
\item \emph{Across a seam:} both sides lie on the same surface $S_f$, and in the periodic parameter domain the two
triangles are adjacent, so the same argument applies; there is no crease.
\item \emph{Along a non-seam edge:} the two patches meet along $\tilde\gamma_e$ at a dihedral angle in $(0,\pi)$ (T), so
the tangent cone at each interior point of the edge is a union of two distinct half-planes.
\item \emph{At a vertex $c_v$:} by T the tangent cone is an embedded disk made of sectors with positive opening angles,
meeting only along common rays; by R the patch corners have positive angles.
\item Hence $h_0$ is locally bi-Lipschitz everywhere, and so bi-Lipschitz.
\end{itemize}

\subsection{Secant approximation}
Let $S$ be the secant map on $T$, refined if necessary: linear on each triangle and equal to $h_0$ at the vertices.
\begin{itemize}
\item On each triangle, $h_0$ is $C^1$ up to the boundary, and the triangles are shape-regular. The standard estimate
gives $\operatorname{Lip}(S-h_0)\le C\,\omega(Dh_0,\mathrm{mesh})/\sin\theta_{\min}$ and $|S-h_0|\le C\,\mathrm{mesh}$.
\item In the intrinsic metric, these per-triangle bounds are global.
\item For a fine enough mesh:
  \begin{itemize}
  \item $\operatorname{Lip}(S-h_0)$ is below the bi-Lipschitz constant of $h_0$, so $S$ is injective;
  \item $|S-h_0|<m$, so $S$ maps every closed cell into its zone, and in particular $S(\inter\sigma)\subseteq\Z(\sigma)$.
  \end{itemize}
\end{itemize}
$S$ is PL, so it is a realization. $\square$

\emph{Standard facts used without separate checking against a source:} existence of shape-regular curved
triangulations of a planar region bounded by $C^2$ curves meeting at positive angles; the per-triangle
secant-approximation estimate; that a map whose tangent cones are embedded disks with positive opening angles is
locally bi-Lipschitz; bi-Lipschitz equivalence of the intrinsic and extrinsic metrics on a compact polyhedron; and the
curved-triangle transformations cited above.

\section{Formal verification in Lean}\label{sec:lean}

The following are proved in Lean~4 with Mathlib (repository directory \texttt{lean/}). The build has no
\texttt{sorry}, and the results depend only on the axioms \texttt{propext}, \texttt{Classical.choice} and
\texttt{Quot.sound}.

\begin{center}\small
\begin{tabular}{p{6.6cm}p{6.4cm}}
\toprule
Lean statement & Content \\
\midrule
\texttt{inner\_eq\_zero\_of\_conical} & at a nearest point $x$, $\langle p-x,x-a\rangle=0$ (Lemma~\ref{lem:cone}) \\
\texttt{inner\_sub\_eq\_norm\_sq\_of\_conical} & $\langle p-x,p-a\rangle=|p-x|^2$ \\
\texttt{not\_isCritical\_of\_conical} & $p$ is not critical (Lemma~\ref{lem:cone}) \\
\texttt{not\_isCritical\_of\_conicalCover} & pointwise form of Corollary~\ref{cor:lambda} \\
\texttt{carrier\_mono}, \texttt{carrier\_eq\_zone\_of\_mono} & carriers are monotone, and equal the zones when the
zones are augmented \\
\texttt{carried\_of\_realization}, \texttt{image\_closedCell\_subset\_carrier} & realizations are carried \\
\texttt{carried\_straightLine}, \texttt{continuous\_straightLine} & Lemma~\ref{lem:carrier} when the carriers
are convex \\
\bottomrule
\end{tabular}
\end{center}

The statements are for an arbitrary real inner product space. Criticality is formalized as
$0\in\conv\{p-x:x\text{ nearest}\}$.

\emph{How the Lean statements correspond to the paper.}
\begin{itemize}
\item The hypothesis \texttt{LocallyConicalAt M a r} is radial closure of $M$ inside the open ball $B(a,r)$. The
paper's definition of $\rho(a)$ implies it for $r<\rho(a)$; the two are not the same statement.
\item Lean's \texttt{Carried} is pointwise on open cells. By \texttt{carrier\_mono} it agrees with the paper's
closed-cell definition.
\item Positivity of $\lambda(M)$ for polyhedra (Corollary~\ref{cor:lambda}) is proved on paper only.
\end{itemize}

\renewcommand{\refname}{参考文献}
\begin{thebibliography}{99}

\bibitem{APS00} L.-E. Andersson, T.~J. Peters, N.~F. Stewart, Equivalence of topological form for curvilinear
geometric objects, \emph{Int.\ J.\ Comput.\ Geom.\ Appl.} 10(6) (2000) 609--622. \unverified{bibliographic data
verified via Crossref; content not read}

\bibitem{ASZ07} L.-E. Andersson, N.~F. Stewart, M. Zidani, Error analysis for operations in solid modeling in the
presence of uncertainty, \emph{SIAM J.\ Sci.\ Comput.} 29(2) (2007) 811--826.

\bibitem{ChazalLieutier07} F. Chazal, A. Lieutier, Stability and computation of topological invariants of solids
in $\mathbb R^n$, \emph{Discrete Comput.\ Geom.} 37 (2007) 601--617.

\bibitem{dREG22} S. di Rocco, D. Eklund, O. Gäfvert, Sampling and homology via bottlenecks, arXiv:2209.01654.
\unverified{exact title}

\bibitem{Durfee83} A.~H. Durfee, Neighborhoods of algebraic sets, \emph{Trans.\ Amer.\ Math.\ Soc.} 276 (1983).
\unverified{page range}

\bibitem{Federer59} H. Federer, Curvature measures, \emph{Trans.\ Amer.\ Math.\ Soc.} 93(3) (1959) 418--491.

\bibitem{Fu85} J.~H.~G. Fu, Tubular neighborhoods in Euclidean spaces, \emph{Duke Math.\ J.} 52 (1985).
\unverified{whole entry}

\bibitem{Jackson95} D.~J. Jackson, Boundary representation modelling with local tolerances, in \emph{Proc.\ 3rd ACM
Symp.\ Solid Modeling and Applications (SMA '95)}, 1995, 247--254.

\bibitem{OCCT} Open Cascade, \emph{OpenCascade Technology documentation}, version 8.0.1 (modeling data: shape
tolerances; Boolean operations specification).

\bibitem{QiShapiro06} J. Qi, V. Shapiro, $\varepsilon$-topological formulation of tolerant solid modeling,
\emph{Computer-Aided Design} 38(4) (2006) 367--377.

\bibitem{QiShapiro06b} J. Qi, V. Shapiro, Geometric interoperability with epsilon solidity, \emph{J.\ Comput.\ Inf.\
Sci.\ Eng.} 6(3) (2006) 213--220.

\bibitem{SSP01} T. Sakkalis, G. Shen, N.~M. Patrikalakis, Topological and geometric properties of interval solid
models, \emph{Graphical Models} 63(3) (2001) 163--175.

\bibitem{Segal90} M. Segal, Using tolerances to guarantee valid polyhedral modeling results, \emph{Computer
Graphics (Proc.\ SIGGRAPH '90)} 24(4) (1990) 105--114.

\bibitem{Shapiro08} V. Shapiro, Homotopy conditions for tolerant geometric queries, in \emph{Reliable Implementation
of Real Number Algorithms}, LNCS 5045, Springer, 2008, 162--180.

\bibitem{SZ08} N.~F. Stewart, M. Zidani, Transfinite interpolation for well-definition in error analysis in solid
modelling, in \emph{Reliable Implementation of Real Number Algorithms}, LNCS 5045, Springer, 2008, 181--192.

\bibitem{Waldhausen68} F. Waldhausen, On irreducible 3-manifolds which are sufficiently large, \emph{Ann.\ of
Math.} 87 (1968) 56--88.

\bibitem{repo} G. Zhang, \texttt{tolerant-solids} (code, proofs, data), \url{https://github.com/eddiezgq/tolerant-solids}.

\end{thebibliography}

\end{document}
